\paragraph{Examples of normal endomorphisms:}

\begin{enumerate}
\def\labelenumi{\arabic{enumi})}
\item
  unitary automorphisms for \(\Phi\) ( \(\S 6\) , no 2), which are characterized by the relation \(u^{-1} = u^{*}\) ( \(\S 1\) , no 8, cor. of prop. 8);
\item
  endomorphisms \(u\) such that \(u^{*}=u\) ; these are called Hermitian endomorphisms.
\end{enumerate}

For every Hermitian endomorphism \(u\) of E, set \(\Phi_{u}(x,y)=\Phi(u(x),y)\) ; we have

\[
\Phi_ {u} (y, x) = \Phi (u (y), x) = \Phi (y, u (x)) = \overline {{\Phi (u (x) , y)}} = \overline {{\Phi_ {u} (x , y)}}.
\]

which shows that \(\Phi_{u}\) is a Hermitian form on E. Conversely, let \(\Psi\) be a Hermitian form on E; since the map \(s_{\Phi}\) from E to E* associated with \(\Phi\) is bijective, there exists, for every \(x \in \mathbf{E}\) , a unique element \(u(x)\) of E such that \(\Psi(x, y) = \Phi(u(x), y)\) ; one verifies easily that \(u\) is a Hermitian endomorphism of E. Thus \(u \to \Phi_{u}\) is a bijection, called canonical, from the set of Hermitian endomorphisms of E onto the set of Hermitian forms on E.

Suppose that A = K or A = K(i); given a bilinear form Ψ on E, there exists, for every \(x \in E\) , a unique element \(u(x)\) , a unique element of E such that Ψ(x, y) = \(\overline{\Phi(u(x), y)}\) ; one sees immediately that u is a semilinear map (for J) from E into itself. Since \(\overline{\Phi(u(x), y)} = \Phi(y, u(x)) = \overline{\Phi(u^{*}(y), x)}\) , one sees that, for Ψ to be symmetric (resp. alternating), it is necessary and sufficient that one have \(u^{*} = u\) (resp. \(u^{*} = -u\) ).

THEOREM 2. --- Let S be a set of endomorphisms of E (resp. of semilinear maps from E into E when A = K or A = K(i)) stable under the map u → u*. Then, if V is a subspace of E stable under S, its orthogonal V⁰ is stable under S. Moreover, E is a direct sum of subspaces stable under S, minimal among the subspaces ≠ \{0\} and stable under S, and pairwise orthogonal.

Indeed, let V be a subspace of E stable under S; for any \(x \in \mathrm{V}^0\) , \(y \in \mathrm{V}\) and \(u \in \mathrm{S}\) , one has \(u^*(y) \in \mathrm{V}\) , whence \(\Phi(y, u(x)) = \Phi(u^*(y), x) = 0\) (resp. \(\Phi(y, u(x)) = \overline{\Phi(u^*(y), x)} = 0\) ), and consequently \(u(x) \in \mathrm{V}^0\) ; thus \(\mathrm{V}^0\) is stable under S, which proves our first assertion. For the second, we proceed by induction on the dimension. \(n\) of E, the case \(n = 0\) being trivial. For \(n \neq 0\) there exists a subspace \(\mathrm{V} \neq \{0\}\) of E stable under S and minimal, for example a stable subspace \(\neq \{0\}\) of minimal dimension. It then suffices to apply to \(\mathrm{V}^0\) the induction hypothesis, since the adjoint of the restriction of \(u\) to \(\mathrm{V}^0\) (with respect to the restriction of \(\Phi\) ) is identical to the restriction to \(\mathrm{V}^0\) of the adjoint of \(u\) .

COROLLARY 1. --- Suppose that A = K or that A = K(i). Let B be a subalgebra of \(\mathfrak{L}_{\mathbf{A}}(\mathbf{E})\) stable under the map \(u \to u^{*}\) . Then E is a semisimple B-module, and is a direct sum of simple submodules that are pairwise orthogonal. The algebra B is semisimple.

Indeed, since every sub-B-module V of E admits a complement, for example \(V^{0}\) the B-module V is semisimple (chap.~VIII, § 3, no. 3, prop. 7). Since every sub-B-module \(\neq\{0\}\) and minimal of E is simple, E is a direct sum of pairwise orthogonal simple submodules. Finally B is a semisimple algebra, since it admits a semisimple and faithful module E whose contramodule is finitely generated (chap.~VIII, § 5, no.~1, prop.~3).

COROLLARY 2. --- With the hypotheses and notations of cor.~1, suppose moreover that the algebra B is commutative. Then all its elements are (normal) semisimple endomorphisms of E. When A = K (resp. A = K(i)), E is a direct sum of pairwise orthogonal simple B-submodules, which are vector spaces of dimension 1 or 2 (resp. 1) over A.

The first assertion follows from chap.~VIII, § 9, no.~1, prop.~2, since B is semisimple. On the other hand every simple B-module is isomorphic to a B-module of the form B/m, where m is a maximal ideal of B, hence to a commutative field L of finite degree over A; when A = K (resp. A = K(i)), the field L is necessarily isomorphic to K or K(i) (resp. K(i)), since K(i) is algebraically closed (chap.~VI, § 2, no.~6, th.~3).

PROPOSITION 4. --- Let u be a normal endomorphism of E. When A is equal to K(i) or to the field of quaternions over K, there exists an orthonormal basis (for Φ) of E consisting of eigenvectors of u. When A = K, u is semisimple and E is a direct sum of subspaces stable under u, pairwise orthogonal, and of dimension 1 or 2.

Let us first examine the case where A is commutative (A = K or A = K(i)). Then the subalgebra B = A{[}u, u*{]} of ℒ\_A(E) is commutative since u is normal; it is stable under the map v → v* by virtue of formulas (32) and (33) of § 1, no.~8. The assertion relative to the case A = K then follows immediately from cor.~2 of th.~2. When A = K(i), this corollary also shows that E is a direct sum of vector subspaces Ax\_i (i = 1, . . ., n) of dimension 1, pairwise orthogonal and stable under u; if we set \(e_{i} = (\Phi(x_{i}, x_{i}))^{-1/2}x_{i}\) , \((e_{i})\) is the desired orthonormal basis.

When A is the field of quaternions over K, it will likewise suffice for us to prove, by virtue of th.~2, that every minimal element of the set of subspaces ≠ \{0\} of E stable under u and u* is of dimension 1. Now such a subspace V necessarily contains an eigenvector x ≠ 0 of u (*), as one sees by noting that the field of quaternions A contains K(i) as an algebraically closed subfield, and by restricting to K(i) the field of scalars of V. It therefore remains only to show that the eigenvector \(x\) of \(u\) is also an eigenvector of \(u^{*}\) . Set \(u(x) = ax (a \in A)\) ; we then have \(\Phi(u(x), x) = a\Phi(x, x) = \Phi(x, x)a = \Phi(x, \bar{a}x)\) since \(\Phi(x, x)\) belongs to the center of A, and on the other hand \(\Phi(u(x), x) = \Phi(x, u^*(x))\) ; it follows that we have \(\Phi(x, u^*(x) - \bar{a}x) = 0\) , and we can therefore write \(u^*(x) = \bar{a}x + z\) , where \(z\) is a vector orthogonal to \(x\) . We therefore have

\[
\Phi (u ^ {*} (x), u ^ {*} (x)) = \bar {a} a \Phi (x, x) + \Phi (z, z) = \Phi (u (x), u (x)) + \Phi (z, z).
\]

Now, since \(u\) is normal, we have

\[
\begin{array}{r l} \Phi (u (x), u (x)) & = \Phi (x, u ^ {*} u (x)) = \Phi (x, u u ^ {*} (x)) \\ & = \Phi (u u ^ {*} (x), x) = \Phi (u ^ {*} (x), u ^ {*} (x)). \end{array}
\]

Consequently we have \(\Phi(z, z) = 0\) , which, by hypothesis, implies \(z = 0\) and \(u^{*}(x) = \bar{a}x\) .Q.E.D.

Remark. --- It follows from th. 2 that the eigenspaces corresponding to two distinct eigenvalues of \(u\) are orthogonal.

PROPOSITION 5. --- Let u be a Hermitian endomorphism of E. The eigenvalues of u belong to K, and there exists an orthonormal basis of E consisting of eigenvectors of u.

When A is equal to K(i) or to the field of quaternions over K, it suffices, by virtue of prop. 4, to prove the first assertion; now, if x is a nonzero eigenvector of u and a is the corresponding eigenvalue, we have, by virtue of the hypothesis u = u*,

\[
a \Phi (x, x) = \Phi (u (x), x) = \Phi (x, u (x)) = \Phi (x, x) \bar {a};
\]

since \(\Phi(x, x)\) is a nonzero element of the center of A, it follows that \(a = \bar{a}\) , hence \(a \in K\) . Consequently a Hermitian matrix has all its eigenvalues in K. Suppose now that A is equal to K. The matrix \(M\) of \(u\) with respect to an orthonormal basis of E is then symmetric; it is therefore a Hermitian matrix if one considers it as a matrix over \(K(i)\) . The first part of the proof then shows that this matrix has all its eigenvalues in K, and hence admits, if \(E \neq \{0\}\) , eigenvectors \(\neq 0\) in E. It follows that every stable subspace for \(u\) which is minimal is necessarily of dimension 1, and the conclusion follows immediately, as in prop. 4.

PROPOSITION 6. --- a) Let Ψ be a Hermitian form (resp. symmetric bilinear form if A = K or A = K(i)) on E. There exists a basis of E which is orthonormal for Φ and orthogonal for Ψ.

\begin{enumerate}
\def\labelenumi{\alph{enumi})}
\setcounter{enumi}{1}
\tightlist
\item
  Suppose A = K or A = K(i), and let Ψ be an alternating bilinear form on E. There exists a basis of E which is orthonormal for Φ and with respect to which the matrix of Ψ is of the form
\end{enumerate}

\[
\left( \begin{array}{c c c c} 0 & a _ {1} & 0 & 0 \ldots 0 \\ - a _ {1} & 0 & 0 & 0 \ldots 0 \\ 0 & 0 & 0 & a _ {2} \ldots 0 \\ 0 & 0 & - a _ {2} & 0 \ldots 0 \\ \ldots \ldots \end{array} \right).
\]

where the \(a_{i}\) are \(\geqslant 0\) in K.

When \(\Psi\) is Hermitian, our assertion follows immediately from prop. 5 and the canonical correspondence between Hermitian forms on E and Hermitian endomorphisms for \(\Phi\) . In the other two cases, let \(u\) be the semilinear map from E to itself defined at the beginning of this n° by the formula \(\Psi(x, y) = \overline{\Phi(u(x), y)}\) . Then \(u^2\) is an A-linear map; when \(\Psi\) is symmetric (resp. alternating), one has \(u = u^*\) (resp. \(u = -u^*\) ), hence \(u^2\) is Hermitian; whence also

\[
\begin{array}{r l} \Phi (u ^ {2} (x), x) & = \overline {{\Phi (x , u ^ {2} (x))}} = \Phi (u ^ {*} (x), u (x)) = \Phi (u (x), u (x)) \\ & (\text { resp. } - \Phi (u (x), u (x))) \end{array}
\]

this shows that the Hermitian form \((x, y) \to \Phi(u^{2}(x), y)\) is positive (resp. negative). Then apply th. 2, and let V be a minimal element of the family of subspaces \(\neq\{0\}\) of E stable under u and under \(u^{*}\) . Since \(u^{2}\) is a Hermitian endomorphism, V contains an eigenvector \(x \neq 0\) of \(u^{2}\) ; set \(u^{2}(x) = ax\) , where \(a \in K\) (prop. 5).

When \(\Psi\) is symmetric, the inequality \(\Phi(u^{2}(x), x) \geqslant 0\) shows that one has \(a \geqslant 0\) . Set \(y = a^{1/2}x + u(x)\) ; we have \(u(y) = a^{1/2}u(x) + ax = a^{1/2}y\) , and Ay is stable under \(u\) and \(u^{*} = u\) . If \(y = 0\) , Ax is stable under \(u\) and \(u^{*}\) ; in all cases, V is a subspace of dimension 1, and our conclusion follows immediately from th. 2.

When \(\Psi\) is alternating, one has \(a \leqslant 0\) ; set

\[
y = (- a) ^ {1 / 2} x + u (x), \quad z = (- a) ^ {1 / 2} x - u (x).
\]

\[
\text { One has } u (y) = - (- a) ^ {1 / 2} z, u (z) = (- a) ^ {1 / 2} y \text { and }
\]

\[
\Phi (y, z) = - a \Phi (x, x) - \Phi (u (x), u (x)) = - \Phi (a x, x) + \Phi (u ^ {2} (x), x) = 0.
\]

If \(y = z = 0\) , one has \(a = 0\) and \(u(x) = 0\) , hence \(Ax\) is stable under \(u\) and \(u^{*}\) , \(V\) is of dimension 1, and the matrix of the restriction of \(\Psi\) to \(V\) is zero since \(\Psi\) is alternating. Otherwise, \(y\) and \(z\) are both \(\neq 0\) , they generate \(V\) , and since they are orthogonal, \(V\) is of dimension 2 and the matrix of the restriction of \(\Psi\) to \(V\) is of the form \(\begin{pmatrix} 0 & b \\ -b & 0 \end{pmatrix}\) . Finally, by virtue of the formula \(\Psi(x', y') = \overline{\Phi(u(x'), y')}\) , one has \(\Psi(x', y') = 0\) when \(x'\) and \(y'\) belong to two subspaces stable under \(u\) and orthogonal for \(\Phi\) . This proves the existence of an orthonormal basis of \(E\) (for \(\Phi\) ) with respect to which the matrix of \(\Psi\) has the indicated form. Q.E.D.

Exercises. --- 1) Suppose the hypotheses at the beginning of § 7 are satisfied; let \(\Phi\) be a positive Hermitian form on E.

\begin{enumerate}
\def\labelenumi{\alph{enumi})}
\item
  For us to have \(\Phi(x,x)\Phi(y,y)=\Phi(x,y)\overline{\Phi(x,y)}\) , it is necessary and sufficient that \(x\) and \(y\) be linearly dependent, or that the plane spanned by \(x\) and \(y\) be isotropic.
\item
  We assume that, for all \(x \in E\) , \(\Phi(x, x)\) is a square in K. Show that for any two vectors x, y in E, we have
\end{enumerate}

\[
\sqrt {\Phi (x + y , x + y)} \leqslant \sqrt {\Phi (x , x)} + \sqrt {\Phi (y , y)}.
\]

If \(\Phi\) is nondegenerate, the two sides of this inequality can be equal only if \(\alpha x + \beta y = 0\) , where \(\alpha\) and \(\beta\) are two elements of K, not both zero, and such that \(\alpha \beta \leqslant 0\) .

\begin{enumerate}
\def\labelenumi{\arabic{enumi})}
\setcounter{enumi}{1}
\tightlist
\item
  We assume \(\mathbf{A} = \mathbf{K}\) or \(\mathbf{A} = \mathbf{K}(i)\) . Let \(X\) be a square Hermitian matrix of order \(n\) on \(\mathbf{A}\) such that for every nonempty subset \(\mathbf{H}\) of \([1, n]\) one has \(X_{\mathbf{H},\mathbf{H}} \geqslant 0\) (notations of prop. 3 of no. 1).
\end{enumerate}

\begin{enumerate}
\def\labelenumi{\alph{enumi})}
\item
  Let \(\lambda\) be an element \(>0\) of K; show that the matrix \(X + \lambda I\) is positive nondegenerate (use prop. 3 of no. 1, arguing by induction on \(n\) ).
\item
  Deduce that the Hermitian matrix X is positive.
\end{enumerate}

\begin{enumerate}
\def\labelenumi{\arabic{enumi})}
\setcounter{enumi}{2}
\item
  We assume that \(\mathbf{A} = \mathbf{K}\) or \(\mathbf{A} = \mathbf{K}(i)\) and that E is finite-dimensional. Let \(\Phi_1, \Phi_2\) be two positive Hermitian forms on E, and let \(V = (\alpha_{ij})\) , \(W = (\beta_{ij})\) be the matrices of these forms with respect to the same basis \((e_i)\) of E. Show that the Hermitian form \(\Phi\) whose matrix \((\gamma_{ij})\) with respect to \((e_i)\) is such that \(\gamma_{ij} = \alpha_{ij}\beta_{ij}\) for every pair of indices, is positive; moreover, if \(\Phi_1\) and \(\Phi_2\) are nondegenerate, the same holds for \(\Phi\) . (In the computation of \(\Phi(x, x)\) express the \(\alpha_{ij}\) in terms of the values \(\Phi_1(c_i, c_i)\) for an orthogonal basis \((c_i)\) of E relative to \(\Phi_1\) .)
\item
  Suppose that A = K or A = K(i). Let R be a Hermitian matrix of order n over A; we say that R has signature \((s, t)\) if the Hermitian form on \(A^{n}\) having R as its matrix with respect to the canonical basis, has signature \((s, t)\) . We denote by \(\Delta_{k}\) the principal minor of R obtained by deleting from R the rows and columns with index \textgreater{} k ; we assume that \(\Delta_{s+t} \neq 0\) and that, for no index \(k < s + t\) , \(\Delta_{k}\) and \(\Delta_{k+1}\) are not simultaneously zero (cf.~§ 6, exerc. 1 b)). Show that if \(\Delta_{k} = 0\) for a \(k < s + t\) , \(\Delta_{k-1}\) and \(\Delta_{k+1}\) have opposite signs (method of exerc. 1 a) of § 6) and that the number s - t is equal to
\end{enumerate}

\[
\operatorname{sgn} \Delta_ {1} + \operatorname{sgn} \left(\Delta_ {1} \Delta_ {2}\right) + \dots + \operatorname{sgn} \left(\Delta_ {s + t - 1} \Delta_ {s + t}\right)
\]

(use exercise 2 of § 6).

\begin{enumerate}
\def\labelenumi{\arabic{enumi})}
\setcounter{enumi}{4}
\item
  We assume that A = K or A = K(i); let R, S be two Hermitian matrices over A, with signatures \((s, t)\) and \((s', t')\) respectively (exerc. 4); show that the matrix \(R \otimes S\) is a Hermitian matrix of signature \((ss' + tt', st' + s't)\) .
\item
  Let K be a maximal ordered field, L a finite-dimensional simple algebra over K. If \((s, t)\) is the signature of the symmetric bilinear form \((x, y) \to \mathrm{Tr}_{\mathrm{L/K}}(xy)\) on L, show that we have:
\end{enumerate}

\(s-t=m\) if L is isomorphic to a matrix algebra of order \(m\) over K;

s - t = 0 if L is isomorphic to a matrix algebra over the field K(i);

\(s - t = -2m\) if L is isomorphic to a matrix algebra of order \(m\) over the field of quaternions over K.

\begin{enumerate}
\def\labelenumi{\arabic{enumi})}
\setcounter{enumi}{6}
\tightlist
\item
  Assume that A satisfies the conditions at the beginning of §7 and that E has finite dimension n. Let Φ be a positive nondegenerate Hermitian form on E.
\end{enumerate}

\begin{enumerate}
\def\labelenumi{\alph{enumi})}
\item
  Show that every similarity for \(\Phi\) is written in a unique way as the product of a homothety with ratio \(\geqslant 0\) in K, and of a unitary transformation.
\item
  For every basis \((a_{i})_{1\leqslant i\leqslant n}\) of E, show that there exists one and only one orthonormal basis \((e_{i})_{1\leqslant i\leqslant n}\) of E satisfying the following conditions: \(1^{\circ}\) for all \(m\) such that \(1\leqslant m\leqslant n\) , the subspace spanned by \(a_1,\ldots ,a_m\) is identical to the subspace spanned by \(e_1,\ldots ,e_m\) ; \(2^{\circ}\) we have \(\Phi (a_i,e_i) > 0\) in K, for every index \(i\) (cf. § 6, no. 1, prop. 1).
\item
  Deduce from \(b\) that for every invertible square matrix \(M\) of order \(n\) on A, there exists one and only one pair of matrices \((L, U)\) of order \(n\) such that \(U\) is a unitary matrix, that \(L = (\lambda_{ij})\) has only zeros below its diagonal and diagonal entries \(\lambda_{ii}\) belonging to K and \(>0\) and finally that \(M = LU\) .
\end{enumerate}

Assume A = K; let L be a finite-dimensional Hermitian space over Λ, whose metric form is positive non-degenerate. Let M, N be two subsets of L, and u a bijection from M onto N such that one has \(e(u(a), u(b)) = e(a, b)\) for any points a, b in M, show that there exists a displacement whose restriction to M is equal to u (argue by induction on the dimension of the affine linear variety spanned by M, as in the Gram-Schmidt orthogonalization process). Does the proposition extend to the case where A equals K(i) or the field of quaternions over K?

\begin{enumerate}
\def\labelenumi{\arabic{enumi})}
\setcounter{enumi}{8}
\item
  Assume the conditions of No. 3 are satisfied, and in addition that A is the field of quaternions over K. Let u be a normal endomorphism of E. Show that E is the direct sum of subspaces \(F_{k}\) ( \(1 \leqslant k \leqslant r\) ), pairwise orthogonal, such that in each of the \(F_{k}\) there exists an orthonormal basis ( \(e_{ik}\) ) ( \(1 \leqslant i \leqslant n_{k}\) ) and that one has \(u(e_{ik}) = \lambda_{k}e_{ik}\) for \(1 \leqslant i \leqslant n_{k}\) , two \(\lambda_{k}\) with distinct indices not being transformed into each other by an inner automorphism of A. Moreover, if ( \(F_{k}'\) ) is a second decomposition of E having the same properties, with a system ( \(\lambda_{k}'\) ) of eigenvalues, one has \(F_{k}' = F_{k}\) (up to a permutation of the indices) and \(\lambda_{k}' = \alpha_{k}\lambda_{k}\alpha_{k}^{-1}\) ; the set of eigenvectors of u for the eigenvalue \(\lambda_{k}\) is the subspace over the commutant \(A_{k}\) of \(\lambda_{k}\) in A, generated by the \(e_{ik}\) ( \(1 \leqslant i \leqslant n_{k}\) ).
\item
  Assume the conditions of No. 3 are satisfied, and in addition that A is equal to K(i) or to the field of quaternions over K. Let u be a normal endomorphism of E.
\end{enumerate}

\begin{enumerate}
\def\labelenumi{\alph{enumi})}
\item
  Show that for a vector subspace F of E to be such that \(u(\mathrm{F}) \subset \mathrm{F}\) , it is necessary and sufficient that F be generated by eigenvectors of u; we then have \(u(\mathrm{F}^{0}) \subset \mathrm{F}^{0}\) and \(u^{*}(\mathrm{F}) \subset \mathrm{F}\) .
\item
  For \(u\) to be unitary (resp. for \(u^{*} = u\) , \(u^{*} = -u\) ), it is necessary and sufficient that for every eigenvalue \(\lambda\) of \(u\) one has \(\lambda \bar{\lambda} = 1\) (resp. \(\bar{\lambda} = \lambda\) , \(\bar{\lambda} = -\lambda\) ).
\item
  We say that a Hermitian endomorphism \(u\) of E is positive (resp. positive and nondegenerate) if the Hermitian form \((x, y) \to \Phi(u(x), y)\) canonically corresponding to it is positive (resp. positive and nondegenerate); for this it is necessary and sufficient that all eigenvalues of \(u\) be \(\geqslant 0\) (resp. \(>0\) ).
\item
  Show that for every integer m\textgreater{} 0, there exists a normal endomorphism v of E such that \(v^{m}=u\) . If u is positive Hermitian, there exists a unique positive Hermitian endomorphism v such that \(v^{m}=u\) , and there exists a polynomial \(f\in K[X]\) such that \(v=f(u)\) ; this last result is also valid when A = K.
\end{enumerate}

\begin{enumerate}
\def\labelenumi{\arabic{enumi})}
\setcounter{enumi}{10}
\tightlist
\item
  Assume the conditions of no. 3 are satisfied, and in addition that A = K. Let u be a normal endomorphism of E.
\end{enumerate}

\begin{enumerate}
\def\labelenumi{\alph{enumi})}
\item
  Let V be a minimal element of the set of subspaces of E not reduced to 0, stable under u and \(u^{*}\) ; show that if V has dimension 2, the restriction of u to V is a direct similarity with multiplier \textgreater{} 0.
\item
  Show that every subspace of E stable under u is also stable under \(u^{*}\) . (Let \(E_{0}\) be the vector space obtained from E by extending the field of scalars to \(\mathrm{K}(i)\) ; \(\Phi\) is the restriction to E of a positive nondegenerate Hermitian form on \(E_{0}\) and u the restriction to E of a normal endomorphism \(u_{0}\) of \(E_{0}\) (for this form); moreover, E is the set of \(x \in E_{0}\) invariants under an involutive semilinear bijection j of \(E_{0}\) , and we have \(u_{0}j = ju_{0}\) . Then apply exerc. 10 a).)
\end{enumerate}

\begin{enumerate}
\def\labelenumi{\arabic{enumi})}
\setcounter{enumi}{11}
\item
  Assume the conditions of No. 3 are satisfied, and in addition that A is equal to K(i) or to the field of quaternions over K. Show that for every endomorphism u of E, there exists an orthonormal basis of E with respect to which the matrix of u has only zeros below the diagonal. (Proceed by induction on the dimension of E, considering an eigenvector of u.)
\item
  Let A be a field satisfying the conditions at the beginning of § 7, and let E, F be two finite-dimensional vector spaces over A, \(\Phi\) (resp. \(\Psi\) ) a positive non-degenerate Hermitian form on E (resp. F). Show that, for every linear map \(u\) from E to F, the adjoint \(u^{*}\) of \(u\) (for \(\Phi\) and \(\Psi\) ; cf. § 1, no. 8) is such that \(u^{*}u\) and \(uu^{*}\) are positive Hermitian endomorphisms (exercise 10 c)) of E and F, respectively.
\item
  Assume the conditions of No. 3 are satisfied. Let u be an endomorphism of E, \(h_{1}\) , \(h_{2}\) the two positive Hermitian endomorphisms of E, such that we have \(h_{1}^{2}=u^{*}u\) , \(h_{2}^{2}=uu^{*}\) (Exercises 13 and 10 d)).
\end{enumerate}

\begin{enumerate}
\def\labelenumi{\alph{enumi})}
\item
  Show that there exists a unitary endomorphism \(\nu\) such that \(u = \nu h_{1} = h_{2}\nu\) and in particular that \(h_{1}\) and \(h_{2}\) are similar (note that \(\bar{u}(0) = \bar{h}_{1}^{-1}(0)\) and that if V is the subspace orthogonal to \(\bar{u}(0)\) , one has \(\Phi(u(x), u(x)) = \Phi(h_{1}(x), h_{1}(x))\) for all \(x \in \mathrm{V}\) . For \(\nu\) to be uniquely determined, it is necessary and sufficient that \(u\) be bijective. In order for one to be able to take \(\nu\) commuting with \(h_{1}\) , it is necessary and sufficient that \(u\) be normal.
\item
  From a), deduce that every square matrix M of order n over A can be written as UDV, where U and V are unitary matrices and D is a diagonal matrix whose entries are \(\geqslant 0\) and have as their squares the eigenvalues of MM*.
\end{enumerate}

\begin{enumerate}
\def\labelenumi{\arabic{enumi})}
\setcounter{enumi}{14}
\tightlist
\item
  Assume the conditions of No. 3 are satisfied. Show that every positive Hermitian matrix H over A can be written in the form \(LL^{*}\) , where \(L = (\lambda_{ij})\) has only zeros below its diagonal and diagonal terms belonging to K and \(\geqslant 0\) ; moreover, L is uniquely determined by these conditions when H is invertible (cf.~exerc. 10 d) and 7 c).
\end{enumerate}

¶ 16) We assume that the conditions of no. 3 are satisfied and, in addition, that A = K(i). Let u be an endomorphism of E.

\begin{enumerate}
\def\labelenumi{\alph{enumi})}
\item
  The set of eigenvalues of \(u\) is contained in the set U of values of \(\Phi(x, u(x))\) when \(x\) ranges over the set of elements of E such that \(\Phi(x, x) = 1\) .
\item
  A subset C of A = K(i) is said to be convex if, for every pair of elements \((\xi, \eta) \in \mathrm{C}^{2}\) and every \(\tau \in K\) such that \(0 \leqslant \tau \leqslant 1\) we have \(\tau\xi + (1 - \tau)\eta \in C\) . Show that the set U is convex. (Reduce to the case where n = 2; by writing u in the form \(\nu + iw\) , where \(\nu\) and w are Hermitian, and replacing u if necessary by \(\lambda u\) , where \(\lambda \in A\) and \(\lambda\overline{\lambda} = 1\) show that everything reduces to proving the following property. Let \(f(\xi_{1}, \xi_{2}) = \xi_{1}\overline{\xi}_{1} + \xi_{2}\overline{\xi}_{2}\) , \(g(\xi_{1}, \xi_{2}) = a\xi_{1}\overline{\xi}_{1} + b\xi_{2}\overline{\xi}_{2}\) ( \(a \in K, b \in K\) ), \(h(\xi_{1}, \xi_{2}) = \alpha\xi_{1}\overline{\xi}_{1} + \beta\xi_{1}\overline{\xi}_{2} + \beta\xi_{2}\overline{\xi}_{1} + \gamma\xi_{2}\overline{\xi}_{2}\) ( \(\alpha \in K, \gamma \in K, \beta \in A\) ); let \((\eta_{1}, \eta_{2}) \in A^{2}\) , \((\zeta_{1}, \zeta_{2}) \in A^{2}\) such that \(f(\eta_{1}, \eta_{2}) = f(\zeta_{1}, \zeta_{2}) = 1\) , \(g(\eta_{1}, \eta_{2}) = g(\zeta_{1}, \zeta_{2}) = 1\) , \(h(\eta_{1}, \eta_{2}) > 0\) , \(h(\zeta_{1}, \zeta_{2}) < 0\) ; there then exists \((\theta_{1}, \theta_{2}) \in A^{2}\) such that \(f(\theta_{1}, \theta_{2}) = 1\) , \(g(\theta_{1}, \theta_{2}) = 1\) and \(h(\theta_{1}, \theta_{2}) = 0\) . One will begin by noting that for every pair \((\xi_{1}, \xi_{2})\) , there exist \(\mu \in A\) such that \(\mu\overline{\mu} = 1\) and \(\beta\mu\xi_{1}\overline{\xi}_{2} + \beta\overline{\mu}\xi_{2}\overline{\xi}_{1} = 0\) , which will make it possible to reduce to the case where \(\beta = 0\) ; use prop. 5 of chap.~VI, § 2, no.~5.)
\item
  Show that if \(u\) is normal, U is the smallest convex set containing all the eigenvalues of \(u\) . Give an example where \(u\) is not normal but where U still possesses the preceding property. (Take as eigenvalues of \(u\) the elements \(\pm 1 \pm i\) as simple eigenvalues, and 0 as a double eigenvalue.)
\end{enumerate}

¶ 17) We assume the conditions of no.~3 are satisfied; for every \(\xi\in\mathbf{A}\) , we denote by \(|\xi|\) the element \(\rho\geqslant0\) of K such that \(\rho^{2}=\xi\bar{\xi}\) (absolute value of \(\xi\) ). For every square matrix \(M=(\alpha_{ij})\) over A, we set \(f(M)=\max_{i}|\alpha_{ii}|\) , \(g(M)=\max_{i,j}|\alpha_{ij}|\) , and denote by \(\varphi(M)\) the largest absolute value of the eigenvalues of \(M\) (it will be shown that this definition makes sense when A is the field of quaternions over K).

\begin{enumerate}
\def\labelenumi{\alph{enumi})}
\tightlist
\item
  Let \(A, B, D\) be three square matrices of order \(n\) over A. Show that if \(D\) is diagonal, then
\end{enumerate}

\[
g ^ {2} (A D B ^ {*}) \leqslant f ^ {2} (D) f (A A ^ {*}) f (B B ^ {*})\tag{1}
\]

(use inequality (1) of no. 1). Deduce that

\[
g (A B B ^ {*} A ^ {*}) \leqslant f (A A ^ {*}) \varphi (B B ^ {*})\tag{2}
\]

(apply prop. 5 to the Hermitian matrix \(BB^{*}\) ). Deduce that, for \(m\) arbitrary square matrices \(A_{i}\) ( \(1 \leqslant i \leqslant m\) ) over A, we have

\begin{enumerate}
\def\labelenumi{(\arabic{enumi})}
\setcounter{enumi}{2}
\tightlist
\item
\end{enumerate}

\[
g ^ {2} (A _ {1} A _ {2} \dots A _ {m}) \leqslant f (A _ {1} A _ {1} ^ {*}) \varphi (A _ {2} A _ {2} ^ {*}) \dots \varphi (A _ {m - 1} A _ {m - 1} ^ {*}) f (A _ {m} ^ {*} A _ {m})\tag{4}
\]

\[
\varphi^ {2} (A _ {1} A _ {2} \dots A _ {m}) \leqslant \varphi (A _ {1} A _ {1} ^ {*}) \varphi (A _ {2} A _ {2} ^ {*}) \dots \varphi (A _ {m} A _ {m} ^ {*})\tag{5}
\]

\[
g ^ {2} (A _ {1} A _ {2} \dots A _ {m}) \leqslant \varphi (A _ {1} A _ {1} ^ {*}) \dots \varphi (A _ {m} A _ {m} ^ {*}).
\]

(For (3), reason by induction starting from (2). For (4), use exerc. 12; finally deduce (5) from (3).)

Show that inequality (3) no longer necessarily holds when one replaces \(A_{m}^{*}A_{m}\) by \(A_{m}A_{m}^{*}\) (observe that one can have \(f(A^{*}A)\neq f(A A^{*})\) ), or when one replaces \(\varphi(A_{i}A_{i}^{*})\) by \(f(A_{i}A_{i}^{*})\) for \(2\leqslant i\leqslant m-1\) (take all the \(A_{i}\) equal to the square matrix all of whose elements are 1).

\begin{enumerate}
\def\labelenumi{\alph{enumi})}
\setcounter{enumi}{1}
\tightlist
\item
  Let u be a normal endomorphism of E; if λ is an eigenvalue of u (an element of K(i) when A = K), then \(\lambda\bar{\lambda}\) is an eigenvalue of u*u. For every normal matrix M (matrix of a normal endomorphism of E with respect to an orthonormal basis), we therefore have φ²(M) = φ(MM*). Deduce that one also has g(M) ≤ φ(M), and more generally, if M₁, . . ., Mₘ are normal,
\end{enumerate}

\begin{enumerate}
\def\labelenumi{(\arabic{enumi})}
\setcounter{enumi}{5}
\tightlist
\item
\end{enumerate}

\[
g (M _ {1} \dots M _ {m}) \leqslant \varphi (M _ {1}) \dots \varphi (M _ {m})\tag{7}
\]

\[
\varphi (M _ {1} \dots M _ {m}) \leqslant \varphi (M _ {1}) \dots \varphi (M _ {m})
\]

(use (4) and (5)).

\begin{enumerate}
\def\labelenumi{\alph{enumi})}
\setcounter{enumi}{2}
\tightlist
\item
  Show that if \(H\) is a positive Hermitian matrix over A, we have \(f(H) = g(H) \leqslant \varphi(H)\) (use (3) and b), by writing \(H = AA^{*}\) ).
\end{enumerate}

¶ 18) We assume the conditions of no. 3 are satisfied.

\begin{enumerate}
\def\labelenumi{\alph{enumi})}
\tightlist
\item
  For every endomorphism u of E, let \((e_{i})\) be an orthonormal basis of E consisting of eigenvectors of \(u^{*}u\) , and let \(\rho_{i}\) be the eigenvalue of \(u^{*}u\) corresponding to the vector \(e_{i}\) ; we set \(s(u) = (\sum_{i=1}^{n} \rho_{i})^{1/2}\) (square root of \(\mathrm{Tr}(u^{*}u)\) when A is commutative); we have \(s(u^{*}) = s(u)\) . If u, v are two endomorphisms of E, and U, V their matrices with respect to the same orthonormal basis of E, show that for the matrix \(UV^{*} + VU^{*}\) , with elements in K, one has \((\mathrm{Tr}(UV^{*} + VU^{*}))^{2} \leqslant 4s(u)s(v)\) (note that \((u^{*} + \lambda v^{*})(u + \lambda v)\) is positive Hermitian for every \(\lambda \in K\) ); deduce that \(s(u + v) \leqslant s(u) + s(v)\) . If in addition A is commutative, one has
\end{enumerate}

\[
\left| \operatorname{Tr} (u v) \right| ^ {2} \leqslant s (u) s (v) \quad \text { and } \quad \left| \operatorname{Tr} (u) \right| \leqslant \sqrt {n} s (u)
\]

(write \(u\) as a product using Exerc. 14 b)).

\begin{enumerate}
\def\labelenumi{\alph{enumi})}
\setcounter{enumi}{1}
\tightlist
\item
  We further assume A is commutative (hence equal to K or K(i)). If \(H_{1}, H_{2}\) are two positive Hermitian square matrices, show that one has \(\mathrm{Tr}(H_{1}H_{2}) \geqslant 0\) (reduce to the case where \(H_{2}\) is a diagonal matrix). Deduce that we have (using the notation of exercise 17)
\end{enumerate}

\[
\left| \operatorname{Tr} (H _ {1} H _ {2}) \right| \leqslant \varphi (H _ {1}) \operatorname{Tr} (H _ {2}) \leqslant \operatorname{Tr} (H _ {1}) \operatorname{Tr} (H _ {2}).
\]

From these inequalities, conclude that for any two endomorphisms \(u, \nu\) of E, we then have \(s(u\nu) \leqslant s(u)s(\nu)\) .

\begin{enumerate}
\def\labelenumi{\alph{enumi})}
\setcounter{enumi}{2}
\tightlist
\item
  Assuming A is still commutative, let \(\prod_{i=1}^{n}(x-\lambda_i)\) be the decomposition into linear factors of the characteristic polynomial of an arbitrary endomorphism \(u\) of E; show that one has \(\sum_{i=1}^{n}|\lambda_i|^2\leqslant(s(u))^2\) and that, for the two sides of this inequality to be equal, it is necessary and sufficient that \(u\) is normal (use Exerc. 12).
\end{enumerate}

\begin{enumerate}
\def\labelenumi{\arabic{enumi})}
\setcounter{enumi}{18}
\item
  Assume the conditions of no. 3 are satisfied. Let M be a square matrix of order n over A; show that for every square submatrix N of M (obtained by deleting from M a certain number of rows and the columns with the same indices as those rows), one has (with the notations of exercise 17) \(\varphi^{2}(N) \leqslant \varphi(MM^{*})\) (apply formula (4) of exerc. 17 appropriately). If in particular M is a normal matrix, \(\varphi(N) \leqslant \varphi(M)\) (cf. exercise 17 b)).
\item
  We assume that the conditions of No. 3 are satisfied and, in addition, that A is equal to K or to K(i). Apply the results of Exercises 17 to 19 to the extension of Φ to the p-th exterior powers (§ 1, No. 9) and to the p-th exterior powers of the endomorphisms or matrices under consideration. In particular, show that, if \(\prod_{i=1}^{n}(X-\lambda_i)\) and \(\prod_{i=1}^{n}(X-\rho_i^2)\) are the decompositions into linear factors of the characteristic polynomials of an
\end{enumerate}

endomorphism \(u\) of E and of the endomorphism \(u^*u\) and if we suppose \(|\lambda_i| \geqslant |\lambda_{i+1}|\) and \(\rho_i \geqslant \rho_{i+1} \geqslant 0\) for \(1 \leqslant i \leqslant n-1\) , one has

\[
\mid \lambda_ {1} \lambda_ {2} \dots \lambda_ {h} \mid \leqslant \rho_ {1} \rho_ {2} \dots \rho_ {h}
\]

for \(1 \leqslant h \leqslant n - 1\) and \(|\lambda_1\lambda_2 \ldots \lambda_n| = \rho_1\rho_2 \ldots \rho_n\) .

\begin{enumerate}
\def\labelenumi{\arabic{enumi})}
\setcounter{enumi}{20}
\tightlist
\item
  We suppose the conditions of No. 3 are satisfied. Let u be a Hermitian endomorphism of E and let \(\prod_{i=1}^{n}(X-\lambda_{i})\) be the decomposition into linear factors of its characteristic polynomial; we assume \(\lambda_{i}\geqslant\lambda_{i+1}\) for \(1\leqslant i\leqslant n-1\) .
\end{enumerate}

\begin{enumerate}
\def\labelenumi{\alph{enumi})}
\item
  Show that the largest (resp. the smallest) of the eigenvalues \(\lambda_{i}\) in K is equal to the largest (resp. the smallest) of the values of \(\Phi(u(x), x)\) when \(x\) ranges over the set of \(x \in E\) such that \(\Phi(x, x) = 1\) . (Reason directly, or apply Exerc. 16 c) by reducing to the case where A = K(i) ( \(\S\) 3, exercise 4).)
\item
  Let \(\Psi_{\mathrm{v}}\) be the restriction to a vector subspace V of E of the Hermitian form \(\Psi\) associated with \(u\) , and let \(u_{\mathrm{v}}\) be the Hermitian endomorphism of V associated with \(\Psi_{\mathrm{v}}\) . Show that \(\lambda_{k}\) is the smallest of the largest eigenvalues of the \(u_{\mathrm{v}}\) as V ranges over the set of vector subspaces of E of dimension \(n - k + 1\) (using Prop. 5).
\end{enumerate}

\begin{enumerate}
\def\labelenumi{\arabic{enumi})}
\setcounter{enumi}{21}
\tightlist
\item
  We assume that the conditions of no. 3 are satisfied, and in addition that A is equal to K(i) or to the field of quaternions over K. Let u, v be two normal endomorphisms of E; let (E \(_{i}\) ) \(_{1\leqslant i\leqslant r}\) (resp. (F \(_{j}\) ) \(_{1\leqslant j\leqslant s}\) ) be the decomposition of E into a direct sum of pairwise orthogonal subspaces, such that in E \(_{i}\) (resp. F \(_{j}\) ) there is an orthonormal basis formed by eigenvectors of u (resp. v) for the same eigenvalue \(\lambda_{i}\) (resp. \(\mu_{j}\) ), and that \(\lambda_{h}\) and \(\lambda_{i}\) (resp. \(\mu_{j}\) and \(\mu_{k}\) ) are not transformed into one another by an inner automorphism of A if \(h\neq i\) (resp. \(j\neq k\) ) (cf. prop. 4 and exerc. 9). For an endomorphism w of E to satisfy uw = wν, it is necessary and sufficient that for every j ( \(1\leqslant j\leqslant s\) ), w(F \(_{j}\) ) be contained in one of the E \(_{i}\) and that the image under w of every eigenvector of v relative to the eigenvalue \(\mu_{j}\) be an eigenvector of u relative to the eigenvalue \(\mu_{j}\) (which implies in particular that \(\mu_{j}\) and \(\lambda_{i}\) are transformed into one another by an inner automorphism of A). Deduce that if this is so, then one has u*w = wν*.
\end{enumerate}

¶ 23) We assume that the conditions of no. 3 are satisfied. Let u and v be two endomorphisms of E.

\begin{enumerate}
\def\labelenumi{\alph{enumi})}
\item
  We assume that u and uv are normal. For vu to be normal, it is necessary and sufficient that v and \(u^{*}u\) commute. (To see that the condition is necessary, use the relation \(u(vu) = (uv)u\) and exerc. 22; to see that it is sufficient, use exerc. 14 a) and 10 d).
\item
  We assume \(u, v\) and \(uv\) normal; let \(\beta\) be the largest eigenvalue of \(uu^*\) , and let \(F\) be the subspace of \(E\) formed by the eigenvectors of \(uu^*\) relative to the eigenvalue \(\beta\) . Show that \(v(F) \subset F\) . (Note that for every normal endomorphism \(u\) and every \(x \in E\) , one has
\end{enumerate}

\[
\Phi (u (x), u (x)) = \Phi (u ^ {*} (x), u ^ {*} (x)).
\]

Reduce to the case where A is equal to K(i) or to the field of quaternions over K; F then admits a basis formed by eigenvectors for u (and \(u^{*}\) ); note that for such a vector z, we have \(\Phi(u^{*}u\nu(z),\nu(z))=\beta\Phi(\nu(z),\nu(z))\) Deduce from this that \(\nu u\) is normal (argue by induction on the number of distinct eigenvalues of \(u^{*}u\) ). If h (resp. \(h'\) ) is the positive Hermitian endomorphism such that \(h^{2}=uu^{*}\) (resp. \(h'^{2}=\nu\nu^{*}\) ) and if we set \(u=hu_{1}\) , \(\nu=h'\nu_{1}\) , where \(u_{1}\) and \(\nu_{1}\) are unitary, h commuting with \(u_{1}\) and \(h'\) with \(\nu_{1}\) (exerc. 14 a)), show that the pairs \((h,h')\) , \((h,\nu_{1})\) and \((h',u_{1})\) are permutable; and conversely. Deduce that \(u^{m}\nu^{n}\) , \(\nu^{n}u^{m}\) , \(u\nu^{*}\) and \(\nu^{*}u\) are then normal (m and n integers \textgreater0 arbitrary).

¶ 24) Assume the conditions of no. 3 are satisfied. Let Γ be a group of automorphisms of E such that every \(u \in \Gamma\) is normal. Show that there exists a decomposition of E into a direct sum of subspaces \(E_k\) ( \(1 \leqslant k \leqslant r\) ) pairwise orthogonal, such that the restriction to \(E_k\) of every \(u \in \Gamma\) is of the form \(\lambda_k v_k\) , where \(\lambda_k\) is an element \textgreater0 of K and where \(v_k\) is a unitary endomorphism of \(E_k\) (decompose each \(u \in \Gamma\) in the form \(h\nu\) , where \(v\) is unitary, \(h\) positive Hermitian, and \(h^2 = uu^*\) (Exerc. 14a); use Exerc. 23b and apply Cor. 2 of Thm. 2 to the algebra generated by the Hermitian endomorphisms \(h\) corresponding to the \(u \in \Gamma\) ). Deduce that if Γ is finite, the elements \(u \in \Gamma\) are unitary.

\begin{enumerate}
\def\labelenumi{\arabic{enumi})}
\setcounter{enumi}{24}
\tightlist
\item
  Suppose that A = K (maximal ordered field). Let L be an n-dimensional Euclidean space over A, whose metric form is positive non-degenerate.
\end{enumerate}

Let S be a nondegenerate affine quadric in L (§ 6, exerc. 25). If S admits a center a and if one takes a as origin in L, show that there exists an orthonormal basis \((e_{i})\) for L such that, with respect to this basis, S has the equation \(\lambda_{1}\xi_{1}^{2} + \cdots + \lambda_{n}\xi_{n}^{2} = 1\) . Moreover, if two orthonormal bases have this property, the elements \(\lambda_{i} \in K\) which correspond to them are the same up to order.

If S admits no center, show that there exists a point b of S and (taking b as origin) an orthonormal basis for L such that, with respect to this basis, S has equation \(\lambda_{1}\xi_{1}^{2} + \cdots + \lambda_{n-1}\xi_{n-1}^{2} + \xi_{n} = 0\) . (If \(\overline{S}\) is the projective quadric such that \(S = L \cap \overline{S}\) , c the pole (at infinity) with respect to \(\overline{S}\) of the hyperplane at infinity \(H_{0}\) , determine b by the condition that the line passing through b and c be perpendicular to the tangent hyperplane to S at the point b).

¶ 26) a) Let K be a commutative field, S a subset of K such that K is a maximal S-orderable field (chap.~VI, § 2, exerc. 8). Let f be a polynomial in K{[}X{]}, L its splitting field. Show that if, for every (total) order structure on K compatible with its ring structure, L admits an ordered extension structure of K, then necessarily L = K. (Argue by contradiction: proceeding as in exerc. 8 e) of chap.~VI, § 2, reduce to the case where one would have {[}L : K{]} = 2. Conclude by noting that if b ∈ K is not a square, there exists a total order structure on K, compatible with its ring structure, for which b \textless{} 0 (cf.~chap.~VI, § 2, n° 3, lemma of th. 1)).

\begin{enumerate}
\def\labelenumi{\alph{enumi})}
\setcounter{enumi}{1}
\tightlist
\item
  Extend the results of n° 3 (prop. 6 excepted) to the case where K is a maximal S-orderable field (*). (Begin by proving the analogue of prop. 5, using a). Then establish the analogue of cor. 2 of th. 2 for the case where the commutative algebra B consists of Hermitian endomorphisms, using prop. 10 of chap.~VIII, § 9, n° 4. Pass to the case of a normal endomorphism u for A = K(i) by noting that one can then write u = v + iw, where v and w are Hermitian and commute. Finally, if A is the quaternion field over K, E₀ the space E considered as a vector space of dimension 2n over K(i), and if one sets Φ(x, y) = Φ₁(x, y) + Φ₂(x, y)j, where Φ₁ and Φ₂ take their values in K(i), note that if u is normal for Φ, it is also normal for Φ₁.)
\end{enumerate}

¶ 27) Let K be a maximal ordered field, E a vector space of dimension n over K, Φ a non-degenerate symmetric bilinear form on E, of signature (s, t) distinct from (n, 0) and (0, n). Let (ei) be an orthogonal basis for Φ, such that Φ(ei, ei) = 1 for 1 ≤ i ≤ s, Φ(ei, ei) = -1 for s + 1 ≤ i ≤ n.~For every orthogonal transformation u ∈ U(Φ), let

\[
U = \left( \begin{array}{c c} M & N \\ P & Q \end{array} \right)
\]

the matrix of \(u\) with respect to \((e_{i})\) , written in the form of a square array of matrices corresponding to the partition of \([1, n]\) into \([1, s]\) and \([s + 1, n]\) .

\begin{enumerate}
\def\labelenumi{\alph{enumi})}
\tightlist
\item
  Prove the relations
\end{enumerate}

\[
\begin{array}{r l} ^ {t} M. M - ^ {t} P. P & = I _ {s} \\ ^ {t} Q. Q - ^ {t} N. N & = I _ {t} \\ ^ {t} M. N - ^ {t} P. Q & = 0. \end{array}
\]

\begin{enumerate}
\def\labelenumi{\alph{enumi})}
\setcounter{enumi}{1}
\item
  Let R be a matrix over K with t rows and s columns, such that \(I_{s} - ^{t}R \cdot R\) is the matrix of a positive non-degenerate Hermitian form (over \(K^{s}\) ). Show that det \((M + \lambda NR)\) does not change sign for \(-1 \leqslant \lambda \leqslant 1\) in K (show, using a), that \(^{t}(M + \lambda NR)(M + \lambda NR)\) is the matrix of a positive non-degenerate symmetric form).
\item
  Set \(\sigma(u) = \sigma(U) = \operatorname{sgn}(\det M)\) ; show that, for any two elements \(u\) , \(\nu\) of the orthogonal group \(\mathbf{U}(\Phi)\) , one has \(\sigma(u\nu) = \sigma(u)\sigma(\nu)\) (using \(a\) ) and \(b\) ). Deduce that the commutator subgroup of the special orthogonal group \(\mathbf{SU}(\Phi)\) is distinct from \(\mathbf{SU}(\Phi)\) (cf.~§ 10, exerc. 9).
\end{enumerate}

\section{Types of quadratic forms}

In this paragraph we assume that A is a commutative field.

\subsection{Types of quadratic forms.}

Given a quadratic form Q (§ 3, n° 4) on a vector space E over A, we shall say that E is the space of definition of Q and that dim(E) is the dimension of Q. Given two quadratic forms Q, Q' on vector spaces E, E' over A, we shall denote by Q τ Q' their direct sum (§ 3, n° 4). Recall that the direct sum of two neutral forms is neutral (§ 4, n° 2).

Introduce the following relation:

“Q and Q' are non-degenerate quadratic forms of finite dimensions over A, and there exist neutral quadratic forms, N, N' such that Q τ N is equivalent to Q' τ N'”.

This relation, which we shall denote by \(Q \sim Q'\) , is manifestly reflexive and symmetric. It is also transitive: indeed, if \(Q\) , \(Q'\) , \(Q''\) are quadratic forms such that \(Q \sim Q'\) and \(Q' \sim Q''\) , there exist neutral quadratic forms \(M\) , \(M'\) , \(N\) , \(N'\) such that \(Q \tau M\) is equivalent to \(Q' \tau M'\) and \(Q' \tau N\) is equivalent to \(Q'' \tau N'\) ; then \(Q \tau (M \tau N)\) is equivalent to \((Q \tau M) \tau N\) , hence to \((Q' \tau M') \tau N\) , and also to \((Q' \tau N) \tau M'\) , hence again to \((Q'' \tau N') \tau M'\) and to \(Q'' \tau (N' \tau M')\) ; since \(M \tau N\) and \(N' \tau M'\) are neutral, one indeed has \(Q \sim Q''\) . The relation \(Q \sim Q'\) is therefore an equivalence relation between \(Q\) and \(Q'\) . It is clear that, if \(Q\) and \(Q'\) are two non-degenerate quadratic forms of finite dimensions and equivalent, one has \(Q \sim Q'\) .

For every quadratic form Q on A, non-degenerate and of finite dimension, we set

\[
\theta (Q) = \tau_ {x} (X \sim Q),\tag{1}
\]

and we shall say that \(\theta(Q)\) is the type of \(Q\) . If \(Q\) and \(Q'\) are two quadratic forms on \(A\) , non-degenerate and of finite dimensions, the relations \(Q \sim Q'\) and \(\theta(Q) = \theta(Q')\) are equivalent.

PROPOSITION 1. --- Let Q and Q' be two nondegenerate quadratic forms on A, of finite dimension. For Q and Q' to be equivalent, it is necessary and sufficient that they have the same dimension and the same type.

The condition is obviously necessary. Suppose it is satisfied. Then there exist neutral forms N, N′ such that Q ⊥ N and Q′ ⊥ N′ are equivalent. Since these two forms have the same dimension, so do N and N′, which are consequently equivalent (§ 4, no. 2, cor. 2 of prop. 2). Hence Q and Q′ are equivalent by virtue of Witt's theorem (§ 4, no. 3, cor. 1 of th. 1).

PROPOSITION 2. --- The relation “there exists a nondegenerate quadratic form Q on A, of finite dimension, such that X = θ(Q)” is collectivizing in X (Set Theory, Chap. II, § 1, no. 4).

Let indeed V be a vector space of infinite dimension over A, S the set of nondegenerate quadratic forms defined on the finite-dimensional subspaces of V, and W the set of \(\theta(Q)\) for \(Q \in \mathfrak{S}\) . It is clear that every nondegenerate quadratic form \(Q'\) of finite dimension over A is equivalent to at least one element of \(\mathfrak{S}\) ; whence \(\theta(Q') \in \mathfrak{W}\) , which proves our assertion.

\subsection{Group of types of quadratic forms.}

We shall equip the set \(\mathfrak{W}\) of types of nondegenerate quadratic forms of finite dimensions over A with a structure of a commutative group. We will define an addition in \(\mathfrak{W}\) by the formula

\[
\mathrm{T} + \mathrm{T} ^ {\prime} = \theta (\mathrm{T} _ {\tau} \mathrm{T} ^ {\prime})\tag{2}
\]

\[
(\mathrm{T}, \mathrm{T} ^ {\prime} \text {   in   } \mathfrak {W}),
\]

This addition is commutative since \(T^{\prime} \tau T\) is equivalent to \(T \tau T^{\prime}\) . It is associative because, if T, \(T^{\prime}\) and \(T^{\prime\prime}\) are elements of W, we have

\[
\begin{array}{r l} & (\mathrm{T} + \mathrm{T} ^ {\prime}) + \mathrm{T} ^ {\prime \prime} \sim (\mathrm{T} + \mathrm{T} ^ {\prime}) \tau \mathrm{T} ^ {\prime \prime} \sim (\mathrm{T} \tau \mathrm{T} ^ {\prime}) \tau \mathrm{T} ^ {\prime \prime} \\ & \qquad \sim \mathrm{T} \tau (\mathrm{T} ^ {\prime} \tau \mathrm{T} ^ {\prime \prime}) \sim \mathrm{T} \tau (\mathrm{T} ^ {\prime} + \mathrm{T} ^ {\prime \prime}) \sim \mathrm{T} + (\mathrm{T} ^ {\prime} + \mathrm{T} ^ {\prime \prime}), \end{array}
\]

whence \((\mathrm{T} + \mathrm{T}' ) + \mathrm{T}'' = \mathrm{T} + (\mathrm{T}' + \mathrm{T}'' )\) since two elements of W that have the same type are equal. Moreover, the addition just defined has an identity element: indeed, it is clear that all neutral forms have the same type \(T_{0}\) , namely that of the zero form of dimension zero; one sees immediately that \(T_{0}\) is the sought neutral element. Finally, the existence, for every \(T \in W\) , of an element opposite to T follows immediately from the following proposition:

PROPOSITION 3. — Let Q be a nondegenerate quadratic form of finite dimension on a vector space V over A. Denote by −Q the quadratic form on V defined by \((-Q)(x) = -Q(x) (x \in V)\) . Then the form \(Q \tau (-Q)\) is neutral.

Indeed, the restriction of \(Q \tau (-Q)\) to the diagonal \(D\) of \(V \times V\) is zero. The index of this form is therefore \(\geqslant\frac{1}{2}\dim(V\times V)\) ( \(\S\) 4, no 2, def. 2), and consequently is equal to \(\frac{1}{2}\dim(V\times V)\) (ibid., formula (4)). It follows that \(Q\tau(-Q)\) is neutral (ibid.).

This allows us to state the following definition:

DEFINITION 1. --- The set of types of nondegenerate quadratic forms of finite dimension over A, equipped with the addition defined by (2), is called the group of types of quadratic forms, or Witt group, of A.

Remarks. --- 1) Every nondegenerate quadratic form Q of finite dimension whose type is zero (that is, such that \(\theta(Q) = T_0\) with the notations above) is a neutral form. Indeed, there exist neutral forms N, N′ such that \(Q \tau N\) is equivalent to N′. This shows that Q is of even dimension, hence that there exists a neutral form \(N_1\) of the same dimension as Q. Since Q and \(N_1\) have the same type, it follows from prop. 1 that they are equivalent, hence that Q is neutral.

\begin{enumerate}
\def\labelenumi{\arabic{enumi})}
\setcounter{enumi}{1}
\tightlist
\item
  For every quadratic form Q of finite dimension over A, denote by \(\delta(Q)\) the class modulo 2 of the dimension of Q. We have
\end{enumerate}

\[
\delta (Q \tau Q ^ {\prime}) = \delta (Q) + \delta (Q ^ {\prime}).
\]

Since every neutral form N has even dimension, we have \(\delta(N) = 0\) ; the relation \(Q \sim Q'\) therefore implies \(\delta(Q) = \delta(Q')\) . Thus the restriction of \(\delta\) to the group \(\mathfrak{W}\) of types of quadratic forms on A is a homomorphism from \(\mathfrak{W}\) into the group \(\mathbf{Z}/(2)\) . This homomorphism is surjective when A has characteristic \(\neq 2\) , but is not so if A has characteristic 2, because a quadratic form of odd dimension is then degenerate since its associated bilinear form is alternating (cf.~§ 5).

\begin{enumerate}
\def\labelenumi{\arabic{enumi})}
\setcounter{enumi}{2}
\tightlist
\item
  Let \(a\) be an element \(\neq 0\) of A. If N is a neutral form, then so is \(aN\) . It follows that the relation \(Q \sim Q'\) implies \(aQ \sim aQ'\) . For every element T of the group \(\mathfrak{W}\) , we set
\end{enumerate}

\[
a. \mathrm{T} = \theta (a \mathrm{T}).\tag{3}
\]

We thus obtain an external law of composition between the group A* of nonzero elements of A and the group W. The following formulas follow immediately from the definition:

\[
a. (\mathrm{T} + \mathrm{T} ^ {\prime}) = a. \mathrm{T} + a. \mathrm{T} ^ {\prime}, \quad a b. \mathrm{T} = a. (b. \mathrm{T})\tag{4}
\]

\[
(a, b \text {   in   } A ^ {*}, T, T ^ {\prime} \text {   in   } \mathfrak {W}).
\]

On the other hand, if \(a\) , \(b\) and \(a + b\) are in \(\mathbf{A}^*\) one does not in general have \((a + b)\) . \(\mathrm{T} = a\) . \(\mathrm{T} + b\) . \(\mathrm{T}\) ( \(\mathrm{T} \in \mathfrak{W}\) ).

PROPOSITION 4. --- Let Q be a nondegenerate quadratic form on a finite-dimensional vector space E over A. Suppose A has characteristic \(\neq 2\) , and let \((x_{1},\ldots,x_{n})\) be an orthogonal basis of V. Let \(T_{1}\) denote the type of the quadratic form \(Q_{1}\) defined on the vector space A and such that \(Q_{1}(1)=1\) . The type of Q is then \(\sum_{i=1}^{n}\mathrm{Q}(x_{i}).T_{1}\) .

Indeed, the form Q is equivalent to

\[
(\mathrm{Q} (x _ {1}) \mathrm{Q} _ {1}) \tau \dots \tau (\mathrm{Q} (x _ {n}) \mathrm{Q} _ {1})
\]

COROLLARY. --- With the hypotheses and notations being those of prop. 3, the elements \(a\) . \(\mathrm{T}_{1}\) ( \(a \in \mathrm{A}^{*}\) ) form a set of generators of the group of types of quadratic forms on A.

Thus, determining the structure of the group of types of quadratic forms on A amounts to finding the Z-linear relations that exist among the elements of the form \(a.T_{1}\) . If \(b \in A^{*}\) , the form \(Q_{1}\) defined in Prop. 4 is clearly equivalent to \(b^{2}Q_{1}\) ; one therefore has \(a.T_{1} = ab^{2}.T_{1}\) , which shows that \(a.T_{1}\) depends only on the class of a modulo the subgroup \((A^{*})^{2}\) of squares of elements of \(A^{*}\) . Moreover, it follows from prop. 3 that one has \((-a).T_{1} = -a.T_{1}\) . However, there exist in general other Z-linear relations among the \(a.T_{1}\) besides those which are deduced from the relations we have just indicated.

PROPOSITION 5. --- Suppose that A is a maximal ordered field. Let Q be a nondegenerate quadratic form of finite dimension over A, and let (s, t) be its signature (§ 7, no. 2, def. 2). Then the type of Q is (s - t)·T₁, and the group W of types of quadratic forms over A is an infinite cyclic group generated by T₁.

Indeed, as \(\mathbf{A}^{*}/(\mathbf{A}^{*})^{2}\) has order 2 and that \((-1)\) . \(\mathrm{T}_{1} = -\mathrm{T}_{1}\) , \(\mathfrak{W}\) is generated by \(\mathrm{T}_{1}\) and is consequently monogenic. For all \(n > 0\) , \(n\) . \(\mathrm{T}_{1}\) is the type of positive nondegenerate quadratic forms of dimension \(n\) ; since these forms are not neutral, we have \(n\) . \(\mathrm{T}_{1} \neq 0\) , which shows that \(\mathfrak{W}\) is infinite. Finally, a form of signature \((s, t)\) is isomorphic, with the notations of Prop. 4, to the direct sum of s forms \(Q_{1}\) and of t forms \(-Q_{1}\) (§ 7, no. 2, th. 1); it follows that its type is \((s - t)\) . \(T_{1}\) .

\subsection{Ring of types of quadratic forms.}

We shall assume, in this section, that A is a field of characteristic ≠ 2.

Given two quadratic forms Q, Q' on vector spaces V, V' over A, we shall call the tensor product of Q and Q', and denote by Q ⊗ Q', the quadratic form on V ⊗ V' whose associated bilinear form is the tensor product (§ 1, no. 9, def. 11) of the bilinear forms associated with Q and Q'. One sees readily that Q ⊗ Q' satisfies the relation

\[
(Q \otimes Q ^ {\prime}) (x \otimes x ^ {\prime}) = Q (x) Q ^ {\prime} (x ^ {\prime}) \quad (x \in V, x ^ {\prime} \in V ^ {\prime}).\tag{5}
\]

If Q and Q' are nondegenerate and of finite dimension, then so is Q ⊗ Q' (§ 1, no. 9, prop. 9).

Let Q, Q′, Q″ be quadratic forms on the vector spaces V, V′, V″. Using the canonical isomorphism from V ⊗ V' onto V' ⊗ V (resp. from (V ⊗ V′) ⊗ V″ onto V ⊗ (V′ ⊗ V″), from (V × V′) ⊗ V″ onto (V ⊗ V″) × (V′ ⊗ V″)), we see immediately that Q ⊗ Q' is equivalent to Q' ⊗ Q (resp. (Q ⊗ Q′) ⊗ Q″ to Q ⊗ (Q′ ⊗ Q″), (Q τ Q′) ⊗ Q″ to (Q ⊗ Q″) τ (Q′ ⊗ Q″)).

Let Q and Q' be two non-degenerate quadratic forms of finite dimensions. If Q is neutral, then so is Q⊗Q'. Indeed, let V, V' be the spaces of definition of Q, Q', with dimensions 2n and n', and let W be a totally singular subspace of V of dimension n (§ 4, no. 2); then W⊗V' is a totally singular subspace, and its dimension is half that of V⊗V'; from this one deduces, as in prop. 3, that Q⊗Q' is neutral. Likewise, Q⊗Q' is neutral whenever Q' is neutral.

From this we deduce that, if Q, Q', Q₁, Q₁' are nondegenerate quadratic forms of finite dimension over A, and if one assumes that θ(Q₁) = θ(Q) and θ(Q₁') = θ(Q'), then one has θ(Q₁ ⊗ Q₁') = θ(Q ⊗ Q'). Indeed, it suffices to verify this in the case where \(Q_{1} = Q \tau N\) and \(Q_{1}' = Q' \tau N'\) , N and \(N'\) being neutral forms; in this case \(Q_{1} \otimes Q_{1}'\) is equivalent to

\[
(Q \otimes Q ^ {\prime}) \tau (Q \otimes N ^ {\prime} \tau Q ^ {\prime} \otimes N \tau N \otimes N ^ {\prime})
\]

and the second parenthesis denotes a neutral form; this proves our assertion.

Now let \(\mathfrak{W}\) be the group of types of quadratic forms over A. Let us define, on the set \(\mathfrak{W}\) a second law of composition, denoted multiplicatively, by the formula

\[
\mathrm{TT} ^ {\prime} = \theta (\mathrm{T} \otimes \mathrm{T} ^ {\prime})\tag{6}
\]

\[
(\mathbf {T}, \mathbf {T} ^ {\prime} \text {   in   } \mathfrak {W}).
\]

It follows immediately from what we have just seen that this composition law is commutative, associative, and distributive with respect to addition. It admits an identity element, namely the type \(T_{1}\) of the quadratic form \(Q_{1}\) defined on the vector space A and such that \(Q_{1}(1)=1\) : indeed, by (5), \(Q_{1}\otimes Q=Q\) for every quadratic form Q. The additive group W, equipped with the multiplication just defined, is therefore a commutative ring with identity element; it is called the ring of types of quadratic forms of A (or the Witt ring of A, when no confusion is to be feared).

Remarks. --- 1) It is clear that, if a is an element of A*, then

\[
a. (\mathrm{TT} ^ {\prime}) = (a. \mathrm{T}) \mathrm{T} ^ {\prime} = \mathrm{T} (a. \mathrm{T} ^ {\prime})\tag{7}
\]

for any elements T, T' of W. One will note moreover that one has a.T = TaT, denoting by Ta the type of the quadratic form aQ1 on A.

\begin{enumerate}
\def\labelenumi{\arabic{enumi})}
\setcounter{enumi}{1}
\tightlist
\item
  Since A has characteristic \(\neq 2\) , every element T of \(\mathfrak{W}\) takes the form \(\sum_{i=1}^{n} a_i \cdot T_1\) where \(a_i \in A^*\) (no. 2, prop. 4). We have
\end{enumerate}

\[
(\sum_ {i = 1} ^ {n} a _ {i}. \mathrm{T} _ {1}) (\sum_ {j = 1} ^ {q} b _ {j}. \mathrm{T} _ {1}) = \sum_ {i, j} a _ {i} b _ {j}. \mathrm{T} _ {1} \quad (a _ {i}, b _ {j} \text {   in   } \mathrm{A} ^ {*}).\tag{8}
\]

\begin{enumerate}
\def\labelenumi{\arabic{enumi})}
\setcounter{enumi}{2}
\tightlist
\item
  Suppose that A is a maximal ordered field. Then the ring W is isomorphic to Z (prop. 5), the integer corresponding to the type of a form of signature \((s, t)\) being s - t (ibid.). Since the tensor product of two forms Q, \(Q'\) of signatures \((s, t)\) , \((s', t')\) is a form of dimension \((s + t)(s' + t')\) , it follows, by an elementary calculation, that the signature of \(Q \otimes Q'\) is \((ss' + tt', st' + ts')\) .
\end{enumerate}

\section{Clifford algebras}

In this paragraph, we shall assume the ring A to be commutative. We shall denote by Q a quadratic form on the A-module E, and by \(\Phi\) the associated bilinear form (\(\S\) 3, no. 4).

\subsection{Definition and universal property of the Clifford algebra.}

DEFINITION 1. --- One calls the Clifford algebra of Q and denotes by C(Q) the quotient algebra of the tensor algebra T(E) of the module E by the two-sided ideal (denoted I(Q)) generated by the elements of the form \(x \otimes x - Q(x)\) . 1 ( \(x \in E\) ).

We will denote by \(\rho_{\mathbb{Q}}\) (or simply \(\rho\) when no confusion is to be feared) the mapping from E into C(Q) composed of the canonical mapping from E into T(E) and the canonical mapping \(\sigma\) of T(E) onto C(Q); the map \(\rho_{\mathbb{Q}}\) is said to be canonical. Note that C(Q) is generated by \(\rho_{\mathbb{Q}}(E)\) , and that, for \(x \in E\) , one has

\[
\rho (x) ^ {2} = \mathrm{Q} (x). 1;\tag{1}
\]

whence, replacing \(x\) by \(x + y\) ( \(x, y\) in E)

\[
\rho (x) \rho (y) + \rho (y) \rho (x) = \Phi (x, y). 1\tag{2}
\]

Example. --- If E admits a basis consisting of a single element \(e\) , T(E) is isomorphic to the polynomial algebra A{[}X{]}, and C(Q) is a quadratic extension of A, with basis \((1, u)\) , where \(u\) is the element \(u = \rho(e)\) and satisfies \(u^2 = Q(e)\) .

Denote by \(T^{h}\) the h-th tensor power \(\otimes E\) in \(\mathrm{T}(\mathrm{E})\) , and let \(T^{+}\) (resp. \(T^{-}\) ) be the sum of \(T^{h}\) for h even (resp. odd).

Since T(E) is the direct sum of \(T^{+}\) and \(T^{-}\) and since I(Q) is generated by elements of \(T^{+}\) , I(Q) is the direct sum of \(T^{+} \cap I(Q)\) and \(T^{-} \cap I(Q)\) , and C(Q) is the direct sum of the two submodules \(C^{+}(Q) = \sigma(T^{+})\) and \(C^{-}(Q) = \sigma(T^{-})\) (also denoted \(C^{+}\) and \(C^{-}\) ). The elements of \(C^{+}\) will be called even (resp. odd). We have the relations

\[
\mathrm{C} ^ {+} \mathrm{C} ^ {+} \subset \mathrm{C} ^ {+}, \quad \mathrm{C} ^ {+} \mathrm{C} ^ {-} \subset \mathrm{C} ^ {-}, \quad \mathrm{C} ^ {-} \mathrm{C} ^ {+} \subset \mathrm{C} ^ {-}, \quad \mathrm{C} ^ {-} \mathrm{C} ^ {-} \subset \mathrm{C} ^ {+}.\tag{3}
\]

In particular \(C^{+}\) is a subalgebra of \(C(Q)\) .

PROPOSITION 1. --- Let f be a linear map of E into an algebra D over A such that \(f(x)^{2} = \mathrm{Q}(x)\) . 1 for every \(x \in \mathbf{E}\) . There exists one and only one homomorphism \(\bar{f}\) of C(Q) into D such that \(f = \bar{f} \circ \rho_{\mathbb{Q}}\) .

The uniqueness of \(\bar{f}\) follows from the fact that C(Q) is generated by \(\rho_{\mathbb{Q}}(\mathrm{E})\) . Let \(h\) be the unique homomorphism of T(E) into D extending \(f\) ( \(h\) is defined by \(h(x_1 \otimes \cdots \otimes x_n) = f(x_1) \ldots f(x_n)\) ). We have

\[
h (x \otimes x - \mathrm{Q} (x). 1) = (f (x) ^ {2} - \mathrm{Q} (x)). 1 = 0,
\]

and consequently h vanishes on I(Q) and defines, by passage to the quotient, the homomorphism \(\bar{f}\) sought.

Prop. 1 expresses that C(Q) is the solution of a universal mapping problem (Ens., chap.~IV, § 3, no. 1).

Take in particular for D the opposite algebra of C(Q) and for \(f\) the map \(\rho\) ; prop. 1 implies that there exists one and only one anti-automorphism \(\beta\) of C(Q) whose restriction to \(\rho(E)\) is the identity; it is called the principal anti-automorphism of C(Q). It is clear that \(\beta^{2} = 1\) .

On the other hand, let Q' be a quadratic form on an A-module E' and f a linear map from E to E' such that Q'∘f = Q. We have \(\rho_{\mathbb{Q}'}(f(x))^{2} = \mathrm{Q}'(f(x)).1 = \mathrm{Q}(x).1\) , and consequently there exists one and only one homomorphism C(f) from C(Q) to C(Q') such that C(f)∘ρq = ρq'∘f.~If f is the identity, C(f) is the identity; if Q'' is a quadratic form on an A-module E'' and g a linear map from E' to E'' such that Q''∘g = Q', we have C(g∘f) = C(g)∘C(f). When E' is a submodule of E and f the canonical injection of E' into E (so that Q' is the restriction of Q to E'), we say that C(f) is the canonical homomorphism of C(Q') into C(Q).

Take in particular Q' = Q and for f the map \(x \to -x\) ; we see that there exists one and only one automorphism \(\alpha\) of C(Q) such that \(\alpha \circ \rho = -\rho\) ; it is called the principal automorphism of C(Q). It is clear that \(\alpha^{2} = 1\) , and that the restriction of \(\alpha\) to C⁺ (resp. C⁻) is the identity (resp. the map \(u \to -u\) ).

PROPOSITION 2. --- Let A' be a commutative ring, φ a homomorphism from A to A', Q' the quadratic form on E' = A' ⊗\_A E deduced from Q by extension of scalars (§ 3, n° 4, prop. 3). There exists one and only one isomorphism j from the algebra A' ⊗\_A C(Q) onto C(Q') such that j (1 ⊗ ρ\_Q(x)) = ρ\_Q′(1 ⊗ x) for all x ∈ E.

It suffices to show that the algebra \(C' = A' \otimes C(Q)\) and the map \(1 \otimes_{\rho_{Q}}\) of \(E'\) into \(C'\) form a solution of the same universal problem as \(C(Q')\) and \(\rho_{Q'}\) . Indeed, let \(D'\) be an algebra over \(A'\) and \(f'\) an A'-linear map from \(E'\) in \(D'\) such that \(f'(x')^{2} = Q'(x')\) . 1 for every \(x' \in E'\) . The map \(g : x \to f'(1 \otimes x)\) from E into \(D'\) (considered as an A-module via the homomorphism \(\varphi\) ) is A-linear and we have \(g(x)^{2} = Q'(1 \otimes x)\) . 1 = Q(x). 1 for all \(x \in E\) . There therefore exists one and only one A-homomorphism \(\bar{g}\) from C(Q) to \(D'\) such that \(\bar{g}(\rho_{Q}(x)) = f'(1 \otimes x)\) . Consequently there exists one and only one A'-homomorphism \(\overline{f'}\) from C' to \(D'\) such that \(\overline{f}'(1 \otimes \rho_{Q}(x)) = f'(1 \otimes x)\) for all \(x \in E\) ; by linearity it follows that \(\overline{f}'((1 \otimes \rho_{Q})(x')) = f'(x')\) for all \(x' \in E'\) . QED.

\subsection{Some operations in the tensor algebra.}

In this section we shall denote by \(e_x\) ( \(x \in E\) ) the linear mapping \(u \to x \otimes u\) of the tensor algebra T(E) into itself.

Lemma 1. — Let f be an element of the dual E* of E. There exists one and only one linear mapping \(i_{f}\) of T(E) into itself such that

\begin{enumerate}
\def\labelenumi{(\arabic{enumi})}
\setcounter{enumi}{3}
\tightlist
\item
\end{enumerate}

\[
i _ {f} (1) = 0\tag{5}
\]

\[
i _ {f} \circ e _ {x} + e _ {x} \circ i _ {f} = f (x). I
\]

for all \(x \in \mathbf{E}\)

(where I denotes the identity mapping). The mapping \(f \to i_f\) of E* into \(\mathfrak{U}(\mathrm{T}(\mathrm{E}))\) is linear. One has \(i_f(\mathrm{T}^n) \subset \mathrm{T}^{n-1}\) , \((i_f)^2 = 0\) , and \(i_f \circ i_g + i_g \circ i_f = 0\) for \(f\) , \(g\) in E*. The mapping \(i_f\) is zero on the subalgebra of

T(E) generated by the kernel of f. The ideal I(Q) is stable under \(i_{f}\) ; by passage to the quotient \(i_{f}\) therefore defines a linear mapping (still denoted \(i_{f}\) ) of C(Q) into itself.

Indeed, formula (5) can be written

\[
i _ {f} (x \otimes u) = - x \otimes i _ {f} (u) + f (x). u \quad (x \in \mathrm{E}, u \in \mathrm{T(E)}).\tag{6}
\]

Since (4) completely determines \(i_{f}\) on \(\mathbf{T}^{0}\) , and since (6) determines \(i_{f}\) on \(\mathbf{T}^{n}\) once its values on \(\mathbf{T}^{n-1}\) are known, the uniqueness of \(i_{f}\) is proved. On the other hand, for \(x\in\mathbf{E}\) and \(u\in\mathbf{T}^{n-1}\) , the right-hand side of (6) is bilinear on \(\mathbf{E}\times\mathbf{T}^{n-1}\) ; this proves the existence of \(i_{f}\) by induction on \(n\) (chap.~III, § 1, no. 2) and also proves, by induction on \(n\) , that \(i_{f}(\mathbf{T}^{n})\subset\mathbf{T}^{n-1}\) . If \(f=ag+bh\) ( \(a,b\) in A, \(g,h\) in E*) it is clear that \(ai_{g}+bi_{h}\) satisfies (4) and (5), hence is equal to \(i_{f}\) . We have

\[
(i _ {f}) ^ {2} \circ e _ {x} = - i _ {f} \circ e _ {x} \circ i _ {f} + f (x) i _ {f} = e _ {x} \circ (i _ {f}) ^ {2}
\]

and, since \((i_f)^2(1) = 0\) we deduce by induction on \(n\) that \((i_f)^2\) is zero on \(\mathbf{T}^n\) . Replacing \(f\) by \(f + g\) , the relation \((i_f)^2 = 0\) gives \(i_f \circ i_g + i_g \circ i_f = 0\) . An induction argument analogous to the preceding ones shows that \(i_f\) is zero on the subalgebra generated by the kernel of \(f\) . Finally, (6) implies that the set of elements \(u\) of I(Q) such that \(i_f(u) \in \mathrm{I}(Q)\) is a left ideal of T(E); moreover, if \(u = (x \otimes x - Q(x).1) \otimes v (x \in E, v \in T(E))\) , one has

\[
\begin{array}{r l} i _ {f} (u) & = f (x) x \otimes v - x \otimes i _ {f} (x \otimes v) - \mathrm{Q} (x) i _ {f} (v) \\ & = f (x) x \otimes v - f (x) x \otimes v + x \otimes x \otimes i _ {f} (v) - \mathrm{Q} (x) i _ {f} (v) \\ & = (x \otimes x - \mathrm{Q} (x). 1) \otimes i _ {f} (v); \end{array}
\]

therefore I(Q) is stable under \(i_{f}\) and the last assertion follows. QED.

Let F be a bilinear form on E; in the rest of this paragraph we shall denote by \(i_x^{\mathrm{F}}\) ( \(x \in \mathrm{E}\) ) the map \(i_f\) corresponding to the linear form \(f: y \to \mathrm{F}(x, y)\) on E.

Lemma 2. --- There exists a linear map \(\lambda_{\mathrm{F}}\) of T(E) into itself such that

\begin{enumerate}
\def\labelenumi{(\arabic{enumi})}
\setcounter{enumi}{6}
\tightlist
\item
\end{enumerate}

\[
\lambda_ {\mathrm{F}} (1) = 1\tag{8}
\]

\[
\lambda_ {\mathrm{F}} \circ e _ {x} = (e _ {x} + i _ {x} ^ {\mathrm{F}}) \circ \lambda_ {\mathrm{F}}\tag{\( (x \in E) \) .}
\]

Whatever the \(f \in \mathbf{E}^*\) , one has

\[
\lambda_ {\mathrm{F}} \circ i _ {f} = i _ {f} \circ \lambda_ {\mathrm{F}}.\tag{9}
\]

Indeed, formula (8) is equivalent to

\[
\lambda_ {\mathbb {F}} (x \otimes u) = x \otimes \lambda_ {\mathbb {F}} (u) + i _ {x} ^ {\mathrm{F}} (\lambda_ {\mathbb {F}} (u)) \quad (x \in \mathrm{E}, u \in \mathrm{T} (\mathrm{E})).\tag{10}
\]

Since (7) determines entirely \(\lambda_{\mathbb{F}}\) on \(\mathrm{T}^{0}\) , and since (10) determines \(\lambda_{\mathbb{F}}\) on \(\mathrm{T}^{n}\) once its values on \(\mathrm{T}^{n-1}\) are known, the uniqueness of \(\lambda_{\mathbb{F}}\) is proved. On the other hand, for \(x \in \mathrm{E}\) and \(u \in \mathrm{T}^{n-1}\) the right-hand side of (10) is bilinear on \(\mathrm{E} \times \mathrm{T}^{n-1}\) ; this proves the existence of \(\lambda_{\mathbb{F}}\) by induction on \(n\) . It remains to prove (9), which we shall do by induction; both sides of (9) vanish on \(\mathrm{T}^{0}\) ; suppose (9) holds on \(\sum_{h=0}^{n-1} \mathrm{T}^{h}\) ; we then have, for \(x \in \mathrm{E}\) and \(u \in \mathrm{T}^{n-1}\) :

\[
\begin{array}{r l} (\lambda_ {\mathbb {F}} \circ i _ {f}) (x \otimes u) & = (- \lambda_ {\mathbb {F}} \circ e _ {x} \circ i _ {f} + f (x) \lambda_ {\mathbb {F}}) (u) \\ & = - (e _ {x} + i _ {x} ^ {\mathrm{F}}) \circ \lambda_ {\mathbb {F}} \circ i _ {f} (u) + f (x) \lambda_ {\mathbb {F}} (u) \\ & = - (e _ {x} + i _ {x} ^ {\mathrm{F}}) \circ i _ {f} \circ \lambda_ {\mathbb {F}} (u) + f (x) \lambda_ {\mathbb {F}} (u) \\ & = (i _ {f} \circ e _ {x} \circ \lambda_ {\mathbb {F}} - f (x) \lambda_ {\mathbb {F}} + i _ {f} \circ i _ {x} ^ {\mathrm{F}} \circ \lambda_ {\mathbb {F}} + f (x) \lambda_ {\mathbb {F}}) (u) \\ & = (i _ {f} \circ (e _ {x} + i _ {x} ^ {\mathrm{F}}) \circ \lambda_ {\mathbb {F}}) (u) = (i _ {f} \circ \lambda_ {\mathbb {F}}) (x \otimes u), \end{array}
\]

whence our last assertion.

Lemma 3. --- Let F and G be two bilinear forms on E. Then we have \(\lambda_{\mathrm{F}} \circ \lambda_{\mathrm{G}} = \lambda_{\mathrm{F+G}}\) . For every bilinear form F on E, \(\lambda_{\mathrm{F}}\) is a bijection from T(E) onto itself.

Indeed, \(\lambda_{\mathbb{F}} \circ \lambda_{\mathbb{G}}\) possesses the characteristic properties (7) and (8) of \(\lambda_{\mathbb{F} + \mathbb{G}}: \text{one has } (\lambda_{\mathbb{F}} \circ \lambda_{\mathbb{G}})(1) = 1\) , and

\[
\begin{array}{r l} \lambda_ {\mathrm{F}} \circ \lambda_ {\mathrm{G}} \circ e _ {x} & = \lambda_ {\mathrm{F}} \circ (e _ {x} + i _ {x} ^ {\mathrm{G}}) \circ \lambda_ {\mathrm{G}} = (e _ {x} + i _ {x} ^ {\mathrm{G}} + i _ {x} ^ {\mathrm{F}}) \circ \lambda_ {\mathrm{F}} \circ \lambda_ {\mathrm{G}} \\ & = (e _ {x} + i _ {x} ^ {\mathrm{F+G}}) \circ \lambda_ {\mathrm{F}} \circ \lambda_ {\mathrm{G}}. \end{array}
\]

On the other hand, if F = 0, we have \(i_x^{\mathbb{F}} = 0\) for all \(x \in \mathbb{E}\) and consequently \(\lambda_{\mathbb{F}} = \mathbb{I}\) which implies that, for every F, we have \(\lambda_{\mathbb{P}} \circ \lambda_{-\mathbb{P}} = \lambda_{-\mathbb{F}} \circ \lambda_{\mathbb{F}} = \mathbb{I}\) .

\subsection{Basis of the Clifford algebra.}

PROPOSITION 3. --- Let Q and Q' be two quadratic forms and F a bilinear form on E such that Q'(x) = Q(x) + F(x, x) for all x ∈ E. The map λ\_F maps the ideal I(Q') onto the ideal I(Q), and defines an isomorphism (denoted \(\overline{\lambda_{\mathbf{F}}}\) ) from the A-module C(Q') onto the A-module C(Q).

As \(\lambda_{\mathbb{F}}\) is a bijection whose \(\lambda_{-\mathbb{F}}\) is the inverse bijection (Lemma 3), it suffices to prove the inclusion \(\lambda_{\mathbb{F}}(\mathrm{I}(\mathrm{Q}^{\prime})) \subset \mathrm{I}(\mathrm{Q})\) . Since \(\mathrm{I}(\mathrm{Q})\) is a left ideal stable under \(i_x^{\mathbb{F}}\) (Lemma 1), (8) shows that the set of \(u \in \mathrm{T}(\mathrm{E})\) such that \(\lambda_{\mathbb{F}}(u) \in \mathrm{I}(\mathrm{Q})\) is a left ideal. It therefore suffices to show that, for any \(u \in \mathrm{T}(\mathrm{E})\) and \(x \in \mathrm{E}\) , one has \(\lambda_{\mathbb{F}}(x \otimes x \otimes u - \mathrm{Q}'(x)u) \in \mathrm{I}(\mathrm{Q})\) . Now, by (8) and Lemma 1, we have

\[
\lambda_ {\mathrm{F}} \circ e _ {x} ^ {2} = (e _ {x} + i _ {x} ^ {\mathrm{F}}) ^ {2} \circ \lambda_ {\mathrm{F}} = (e _ {x} ^ {2} + \mathrm{F} (x, x)) \circ \lambda_ {\mathrm{F}},
\]

whence

\[
\begin{array}{r l} \lambda_ {\mathbf {F}} (x \otimes x \otimes u - \mathrm{Q} ^ {\prime} (x) u) & = (e _ {x} ^ {2} + \mathrm{F} (x, x) - \mathrm{Q} ^ {\prime} (x)) \circ \lambda_ {\mathbf {F}} (u) \\ & = (x \otimes x - \mathrm{Q} (x)) \otimes \lambda_ {\mathbf {F}} (u) \in \mathrm{I} (\mathrm{Q}). \end{array}
\]

Lemma 4. — If the quadratic form Q is zero, then C(Q) is none other than the exterior algebra of E.

Indeed, the exterior algebra of E is none other than the quotient of T(E) by the two-sided ideal J generated by the elements of the form \(a \otimes v \otimes a\) ( \(a \in E\) , \(v \in T(E)\) ) (chap.~III, § 5, no.~5 and no.~9). It is clear that I(Q) ⊂ J. It therefore suffices to show that \(a \otimes v \otimes a \in I(Q)\) . This is evident if \(v \in T^0\) . Suppose this assertion has been proved for \(n-1 \leqslant \nu \in \sum_{h=0}^{\infty} \mathrm{T}^{h}\) , and let \(x \in \mathrm{E}\) and \(u \in \mathrm{T}^{n-1}\) ; we have

\[
\begin{array}{c} \boldsymbol {a} \otimes \boldsymbol {x} \otimes \boldsymbol {u} \otimes \boldsymbol {a} = (\boldsymbol {a} + \boldsymbol {x}) \otimes (\boldsymbol {a} + \boldsymbol {x}) \otimes \boldsymbol {u} \otimes \boldsymbol {a} \\ - \boldsymbol {a} \otimes \boldsymbol {a} \otimes \boldsymbol {u} \otimes \boldsymbol {a} - \boldsymbol {x} \otimes \boldsymbol {a} \otimes \boldsymbol {u} \otimes \boldsymbol {a} - \boldsymbol {x} \otimes \boldsymbol {x} \otimes \boldsymbol {u} \otimes \boldsymbol {a} \end{array}
\]

and the four terms on the right-hand side belong to I(Q).

Suppose in particular that the module E admits a basis \((x_{i})_{i\in \mathbf{L}}\) , and let us totally order the index set L. We know (chap.~III, §5, no.~6) that the exterior algebra of E admits as a basis the family consisting of the elements \(x_{\mathrm{H}}\) where H ranges over the set of finite subsets of L and where \(x_{\mathrm{H}}\) is the element \(x_{h_1}\wedge \ldots \wedge x_{h_q}\) , \((h_1,\ldots ,h_q)\) denoting the strictly increasing sequence of elements of H.

On the other hand, consider the bilinear form F defined by \(\mathrm{F}(x_{i}, x_{j}) = -\Phi(x_{i}, x_{j})\) if i \textgreater{} j, \(\mathrm{F}(x_{i}, x_{j}) = 0\) if i \textless{} j and \(\mathrm{F}(x_{i}, x_{i}) = -\mathrm{Q}(x_{i})\) . It is clear that \(\mathrm{Q}(x) + \mathrm{F}(x, x) = 0\) ; with the notations of Prop. 3, it follows from this proposition and Lemma 4 that \(\bar{\lambda}_{\mathrm{F}}\) is an isomorphism of \(\wedge \mathrm{E}\) onto C(Q), which is therefore a free A-module. Let us prove, by induction on the number of elements of H, that we have

\[
\overline {{{{\lambda}}}} _ {\mathrm{F}} (x _ {\mathrm{H}}) = \rho (x _ {h _ {1}}) \dots \rho (x _ {h _ {q}})\tag{11}
\]

\((\mathrm{where}\ (h_{1},\ldots,h_{q})\) is the strictly increasing sequence of elements of H). This is evident if H has 0 or 1 element. Suppose (11) holds for sets having at most \(q-1\) elements. Consider a subset H of \(q\) elements, denote \(j\) its smallest element, and let us write \(\mathrm{H}=\{j\}\cup\mathrm{K}\) where K is a subset with \(q-1\) elements. By (8) and the induction hypothesis, we have

\[
\overline {{{\lambda}}} _ {\mathrm{F}} (x _ {\mathrm{H}}) = \overline {{{\lambda}}} _ {\mathrm{F}} (x _ {j} \wedge x _ {\mathrm{K}}) = \rho (x _ {j}) \overline {{{\lambda}}} _ {\mathrm{F}} (x _ {\mathrm{K}}) + i _ {x _ {j}} ^ {\mathrm{F}} (\overline {{{\lambda}}} _ {\mathrm{F}} (x _ {\mathrm{K}})) = x _ {\mathrm{H}} ^ {\prime} + i _ {x _ {j}} ^ {\mathrm{F}} (x _ {\mathrm{K}} ^ {\prime}),
\]

setting, for every finite subset J of L, \(x_{J}^{\prime}=\rho(x_{j_{1}})\ldots\rho(x_{j_{s}})\) , where \((j_{1},\ldots,j_{s})\) is the increasing sequence of elements of J. Now, for \(i\in\mathbf{K}\) , \(x_{i}\) belongs to the kernel of the linear form \(y\to\mathrm{F}(x_{j},y)\) , hence \(i_{x_{j}}^{\mathrm{F}}(x_{\mathrm{K}}^{\prime})=0\) (Lemma 1). This proves the desired result.

We have therefore proved the following theorem:

THEOREM 1. --- Suppose that the A-module E admits a basis \((x_{i})_{i\in L}\) , the index set L being equipped with a total order structure. For every finite subset H of L, set \(x_{\mathrm{H}} = \rho(x_{h_1})\rho(x_{h_2})\ldots \rho(x_{h_q})\) , where \((h_1,\ldots,h_q)\) is the strictly increasing sequence of the elements of H. Then the elements \(x_{\mathrm{H}}\) form a basis of the A-module C(Q).

COROLLARY 1. --- If E is a free module of dimension n, then C(Q) is a free module of dimension \(2^{n}\) ; moreover, if \(n > 0\) , C \(^{+}\) and C \(^{-}\) are free modules of dimension \(2^{n-1}\) .

This follows immediately from the properties of binomial coefficients.

COROLLARY 2. --- If E is a free module, the canonical mapping \(\rho\) from E into C(Q) and the mapping \(a \to a\) .1 of A into C(Q) are injective.

COROLLARY 3. --- Suppose that E is the direct sum of two free submodules \(E_{1}\) and \(E_{2}\) . Let \(Q_{i}\) be the restriction of Q to \(E_{i}\) and \(p_{i}\) the canonical map from C(Q \(_{i}\) ) into C(Q) (i = 1, 2). Then the

linear map \(p\) of \(\mathrm{C(Q_{1})\otimes C(Q_{2})}\) into \(\mathrm{C(Q)}\) deduced from the bilinear map \((a,b)\to p_{1}(a)p_{2}(b)\) of \(\mathrm{C(Q_{1})\times C(Q_{2})}\) in \(\mathrm{C(Q)}\) is a bijection.

It suffices indeed to consider the basis of E obtained by taking the union of a basis of \(E_{1}\) and a basis of \(E_{2}\) .

COROLLARY 4. --- With the hypotheses and notation of Cor. 3, suppose further that \(E_{1}\) and \(E_{2}\) are orthogonal, and transport to \(C(Q_{1})\otimes C(Q_{2})\) , by means of the bijection p, the algebra structure of \(C(Q)\) . If \(a_{i}\) and \(b_{i}\) are even or odd elements of \(C(Q_{i})\) (i=1, 2), we have \((a_{1}\otimes a_{2})(b_{1}\otimes b_{2})=\varepsilon(a_{1}b_{1})\otimes(a_{2}b_{2})\) , with \(\varepsilon=1\) unless \(a_{2}\) and \(b_{1}\) are odd, in which case \(\varepsilon=-1\) .

It suffices indeed to prove that \(p_2(a_2)p_1(b_1) = \varepsilon p_1(b_1)p_2(a_2)\) , and for that one may assume that \(p_2(a_2)\) (resp. \(p_1(b_1)\) ) is a product \(x_1 \ldots x_h\) (resp. \(y_1 \ldots y_k\) ) of elements of \(\rho_{\mathbb{Q}}(\mathbf{E}_2)\) (resp. \(\rho_{\mathbb{Q}}(\mathbf{E}_1)\) ). Since \(\mathbf{E}_1\) and \(\mathbf{E}_2\) are orthogonal, one has

\[
x _ {i} y _ {j} + y _ {j} x _ {i} = \Phi (x _ {i}, y _ {j}) = 0,
\]

whence

\[
x _ {1} \dots x _ {h} y _ {1} \dots y _ {k} = (- 1) ^ {h k} y _ {1} \dots y _ {k} x _ {1} \dots x _ {h}.
\]

The conclusions of Cor. 3 and 4 remain true if one omits the hypothesis that \(\mathbf{E}_1\) and \(\mathbf{E}_2\) are free modules (cf.~exerc. 1).

\subsection{Structure of the Clifford algebra.}

In this section, we will assume that A is a field, that E is a vector space of finite dimension m over A, and that the quadratic form Q is non-degenerate (which, according to th. 1 of § 5, no. 1, requires that m be even if A has characteristic 2). Since E is free, the canonical map ρ is injective (no. 3, cor. 2 of th. 1). We shall henceforth identify E with its image in C(Q).

THEOREM 2. --- Suppose that the dimension of E is an even number \(m = 2r\) and that Q be neutral ( \(\S 4, \text{n}^{\circ} 2\) ). Then the algebra C(Q) is separable (Chap.~VIII, \(\S 7, \text{n}^{\circ} 5\) , Def. 1) and is isomorphic to the algebra of endomorphisms of a vector space of dimension \(2^{r}\) on A. Moreover, if \(m > 0\) , C \(^{+}\) (Q) is separable and is the direct sum of two ideals isomorphic to the algebra of endomorphisms of a vector space of dimension \(2^{r-1}\) over A.

Indeed, since Q is neutral, one can decompose E as a direct sum of two totally singular subspaces N and P of dimension r (§ 4, no. 2, cor. 1 of prop. 2). Since the restriction of Q to N is zero, the subalgebra S of C(Q) generated by N is identified with the exterior algebra of N (no. 3, lemma 4). For \(n \in \mathbb{N}\) , we shall denote by \(e_n'\) the map \(t \to nt\) from S into itself.

Let \((n_{1},\ldots,n_{r})\) be a basis of N; we shall denote by \((p_{1},\ldots,p_{r})\) the basis of P such that \(\Phi(n_{i},p_{j})=\delta_{ij}\) ( \(\S 4\) , no. 2, prop. 2). For \(p\in P\) , we shall denote by \(p'\) the linear form \(n\to\Phi(n,p)\) on N, and \(i_{p}\) the endomorphism of S obtained by passing to the quotient from the endomorphism \(i_{p'}\) of T(N) associated with \(p'\) as was stated in Lemma 1 of No. 2. By (5), we have

\[
e _ {n} ^ {\prime} \circ i _ {p} + i _ {p} \circ e _ {n} ^ {\prime} = \Phi (n, p) \quad (n \in \mathrm{N}, p \in \mathrm{P}).\tag{12}
\]

Let us set, for \(x = n + p \in \mathbf{E}\) (with \(n \in \mathbf{N}\) and \(p \in \mathbf{P}\) ), \(s(x) = e_n' + i_p\) . It is clear that \(s\) is a linear map from \(\mathbf{E}\) into \(\mathfrak{L}(\mathbf{S})\) . Since we have

\[
s (x) ^ {2} = \left(e _ {n} ^ {\prime} + i _ {p}\right) ^ {2} = Q (n) + \Phi (n, p) = Q (x)
\]

By virtue of (12) and Lemma 1 (no. 2), s extends to a homomorphism (which we shall still denote by s) from C(Q) into ℒ(S) (no. 1, prop. 1). We shall show that this homomorphism is surjective, which, since C(Q) and ℒ(S) are both of dimension 2²ʳ, will imply that s is an isomorphism and will prove our first assertion.

Indeed, let I denote the interval \([1, r]\) . For every subset H of I, we set \(H' = I - H\) and we denote by \(n_{H}\) (resp. \(p_{H}\) ) the product of \(n_{i}\) (resp. \(p_{i}\) ) for \(i \in H\) , arranged in increasing order of indices. Recall that the \(n_{H}\) form a basis of S (no. 3, th. 1). Finally, for any two subsets H, K of I, set \(x_{H,K} = n_{H}p_{I}n_{K'}\) . We shall show that the elements \(s(x_{H,K})\) of \(s(C(Q))\) generate L(S). Now, if \(j \notin H\) , one has \(s(p_{j})(n_{H}) = i_{p_{j}}(n_{H}) = 0\) by lemma 1, since the \(n_{i}\) for \(i \in H\) belong to the kernel of the linear form \(n \to \Phi(n, p_{j})\) on N; on the other hand we have

\[
s \left(p _ {j}\right) \left(n _ {j} n _ {\mathrm{H}}\right) = \left(i _ {p _ {j}} \circ e _ {n _ {j}} ^ {\prime}\right) \left(n _ {\mathrm{H}}\right) = \Phi \left(p _ {j}, n _ {j}\right) n _ {\mathrm{H}} - n _ {j}. s \left(p _ {j}\right) \left(n _ {\mathrm{H}}\right) = n _ {\mathrm{H}}
\]

(by (12)). Since \(s\) is a homomorphism, we deduce, for any two subsets H, K of I, that \(s(p_{\mathrm{K}})(n_{\mathrm{H}}) = 0\) if \(\mathrm{K} \not\subset \mathrm{H}\) , and that \(s(p_{\mathbf{K}}) (n_{\mathbf{H}}) = \pm n_{\mathbf{H} - \mathbf{K}}\) if \(\mathrm{K} \subset \mathrm{H}\) . Since, for \(\mathrm{M} \subset \mathrm{I}\) and \(\mathrm{L} \subset \mathrm{I}\) , we have by definition \(s(n_{\mathbf{M}})(n_{\mathbf{L}}) = n_{\mathbf{M}}n_{\mathbf{L}}\) , and that \(n_{\mathbf{M}}n_{\mathbf{L}}\) is zero if \(\mathrm{M} \cap \mathrm{L} \neq \emptyset\) and is equal to \(\pm n_{\mathbf{M} \cup \mathbf{L}}\) in the contrary case, we conclude from the preceding that, for arbitrary subsets \(\mathrm{H}\) , \(\mathrm{K}\) , \(\mathrm{L}\) of \(\mathrm{I}\) , \(s(x_{\mathbf{H},\mathbf{K}})(n_{\mathbf{L}}) = s(n_{\mathbf{H}})s(p_{\mathbf{I}})s(n_{\mathbf{K}^{\prime}})(n_{\mathbf{L}})\) is zero if \(\mathrm{K} \neq \mathrm{L}\) and is equal to \(\pm n_{\mathbf{H}}\) if \(\mathrm{K} = \mathrm{L}\) . This shows that the \(s(x_{\mathbf{H},\mathbf{K}})\) generate \(\mathfrak{L}(\mathrm{S})\) and completes the proof of the first assertion.

To prove the second assertion, set \(S^{+} = S \cap C^{+}\) and \(S^{-} = S \cap C^{-}\) ; it is clear that \(S^{+}\) (resp. \(S^{-}\) ) is the subspace of S generated by the \(n_{\mathrm{H}}\) such that H has an even (resp. odd) number of elements, that S is the direct sum of \(S^{+}\) and \(S^{-}\) , and that \(s(C^{+})\) leaves \(S^{+}\) and \(S^{-}\) stable. Consequently \(s\) maps \(C^{+}\) into a subalgebra of \(\mathfrak{L}(S)\) , isomorphic to the product \(\mathfrak{L}(S^{+}) \times \mathfrak{L}(S^{-})\) ; the restriction of \(s\) to \(C^{+}\) is an isomorphism of \(C^{+}\) onto this subalgebra, since \(s\) is injective and \(C^{+}\) and \(\mathfrak{L}(S^{+}) \times \mathfrak{L}(S^{-})\) both have dimension \(2^{2r-1}\) (no. 2, cor. 1 of th. 1). QED.

COROLLARY. --- If m is even, but Q of arbitrary index, the algebra C(Q) is a central simple algebra of dimension \(2^{m}\) . Moreover, if m \textgreater{} 0, the subalgebra C⁺(Q) is separable, and its center Z has dimension 2 over A. When Z is a field, Z is a separable quadratic extension of A and C⁺(Q) is simple; otherwise, Z is a direct product of two fields isomorphic to A, and C⁺(Q) is then a direct product of two simple subalgebras of dimensions \(2^{m-2}\) .

Indeed, let A' be the algebraic closure of A, and Q' the quadratic form on \(E' = A' \otimes_{A} E\) obtained from Q by extension of scalars. We have seen that \(C(Q')\) is isomorphic to \(A' \otimes_{A} C(Q)\) (prop. 2), and it is clear that \(C^{+}(Q')\) is isomorphic to \(A' \otimes_{A} C^{+}(Q)\) . Since Q' is neutral (§ 4, no. 2, cor. 2 of prop. 3), the corollary is an immediate consequence of th. 2 and the permanence theorems of chap.~VIII, § 7.

Remarks. --- 1) Since the algebra C(Q) is simple, it has only one class of irreducible representations; these are called the spinor representations; when one fixes one's attention on one of these representations, say \(\tau\) , the elements of the space in which \(\tau\) takes place are called spinors. If Q is neutral, the restriction of \(\tau\) to \(C^{+}(Q)\) is, like that of \(s\) , sum of two inequivalent absolutely irreducible representations; the elements of the subspaces carrying these two representations are called semi-spinors. In the general case, if \(C^{+}(Q)\) is not simple, the restriction of \(\tau\) to \(C^{+}(Q)\) must, since it is faithful, contain subrepresentations belonging to each of the two classes of irreducible representations of \(C^{+}(Q)\) , and hence is a sum of two inequivalent absolutely irreducible representations, since this is the case after extending scalars to the algebraic closure A' of A. On the other hand, if \(C^{+}(Q)\) is simple, it has only one class of irreducible representations, and the restriction of \(\tau\) to \(C^{+}(Q)\) is therefore irreducible, since it decomposes under scalar extension to A' into two non-equivalent representations.

\begin{enumerate}
\def\labelenumi{\arabic{enumi})}
\setcounter{enumi}{1}
\tightlist
\item
  Suppose A has characteristic \(\neq 2\) , and let \((x_{1},\ldots,x_{m})\) ( \(m=2r\) ) be an orthogonal basis of E. Put
\end{enumerate}

\[
z = 2 ^ {r} x _ {1} \dots x _ {m} \in \mathrm{C(Q)};
\]

since \(x_{i}x_{j} + x_{j}x_{i} = 0\) for \(i \neq j\) , one has \(zx_{j} = -x_{j}z\) , which implies that \(z\) belongs to the center \(Z\) of \(C^{+}(Q)\) without belonging to \(A\) . We have

\[
z ^ {2} = 2 ^ {2 r} (- 1) ^ {r} \mathrm{Q} (x _ {1}) \dots \mathrm{Q} (x _ {m}) = (- 1) ^ {r} \mathrm{D}
\]

denoting by D the discriminant of \(\Phi\) with respect to the basis \((x_{j})\) (cf.~exerc. 9).

THEOREM 3. --- Suppose that the dimension of E is an odd number \(m = 2r + 1\) (so that A has characteristic \(\neq 2\) ).

\begin{enumerate}
\def\labelenumi{\alph{enumi})}
\item
  The algebra C \(^{+}\) (Q) is central simple. If Q has maximum index r, C \(^{+}\) (Q) is isomorphic to the algebra of endomorphisms of a vector space of dimension 2 \(^{r}\) over A.
\item
  The algebra C(Q) is separable. Its center Z has dimension 2, and C(Q) is isomorphic to Z⊗ \(_{A}\) C \(^{+}\) (Q), hence is simple or a direct product of two simple subalgebras.
\end{enumerate}

Indeed, let \(x_0\) be a non-isotropic vector of E, and F the orthogonal complement of \(x_0\) ; let us denote by \(Q_1\) the quadratic form \(y \to -Q(x_0)Q(y)\) on F; it is clear that \(Q_1\) is nondegenerate. Since \(x_0y = -yx_0\) (for \(y \in F\) ), we have \((x_0y)^2 = -Q(x_0)Q(y) = Q_1(y)\) , and consequently the map \(y \to x_0y\) of F into \(C^+(Q)\) extends to a homomorphism \(h\) of \(C(Q_1)\) in \(C^+(Q)\) (no. 1, prop. 1). But \(C(Q_1)\) is simple (th. 2) and has the same dimension \(2^{2r}\) as \(C^{+}(Q)\) ; this implies, since \(h(1) = 1\) , that \(h\) is an isomorphism. Moreover, if \(Q\) has index \(r\) , one can choose \(x_0\) in such a way that \(Q_1\) also has index \(r\) ( \(\S 4\) , no. 2, prop. 3), which proves \(a\) ).

Now let \((x_{1},\ldots ,x_{2r})\) be an orthogonal basis of F; set \(z = x_0x_1\dots x_{2r}\) . One immediately verifies that \(z\) commutes with \(x_{j}\) for \(j = 0,\dots,2r\) , hence belongs to the center of C(Q). Let Z be the subspace of C(Q) generated by 1 and \(z\) ; it is a subalgebra of the center of C(Q) and a quadratic extension of A, since \(z\) is odd and \(z^2\) is equal to the scalar \((-1)^r\mathrm{Q}(x_0)\ldots \mathrm{Q}(x_{2r})\) . Consider the homomorphism \(\theta\) of \(\mathbf{Z}\otimes_{\mathbf{A}}\mathbf{C}^{+}(\mathbf{Q})\) into C(Q) defined by \(\theta (u\otimes v) = uv\) . Since \(z\in \mathbb{C}^{-}\) and is invertible, the map \(u\to zu\) is an isomorphism of the module \(\mathbf{C}^+\) onto \(\mathbb{C}^{-}\) , which implies that \(\theta (\mathrm{Z}\otimes \mathrm{C}^{+})\) contains \(\mathbf{C}^{+}\) and \(\mathbb{C}^{-}\) , hence coincides with C(Q). Since \(\mathbf{Z}\otimes \mathbf{C}^{+}\) and C(Q) have the same dimension \(2^{2r + 1},\theta\) is an isomorphism; this proves \(b)\) taking into account the results of Chap.~VIII, § 7.

Note. --- The discriminant D of \(\Phi\) with respect to the basis \((x_{j})_{(j=0,\ldots,2r)}\) is equal to \(2^{2r+1}Q(x_{0})\ldots Q(x_{2r})\) . Consequently, Z is generated by 1 and by the odd element \(z' = 2^{r+1}z\) such that \(z'^{2} = (-1)^{r}2D\) . The algebra C(Q) is therefore simple if and only if \(2(-1)^{r}D\) is not a square in A.
